\documentclass[11pt]{article}
\usepackage[margin=1in]{geometry}
\usepackage{amsmath,amssymb,amsthm,mathtools}
\usepackage[T1]{fontenc}
\usepackage{lmodern}
\usepackage{microtype}
\usepackage[colorlinks=true,linkcolor=blue,urlcolor=blue,citecolor=blue]{hyperref}
\newtheorem{theorem}{Theorem}[section]
\newtheorem{proposition}[theorem]{Proposition}
\newtheorem{lemma}[theorem]{Lemma}
\newtheorem{corollary}[theorem]{Corollary}
\theoremstyle{definition}\newtheorem{definition}[theorem]{Definition}
\theoremstyle{remark}\newtheorem{remark}[theorem]{Remark}
\newcommand{\SLE}{\mathrm{SLE}_6}
\newcommand{\HH}{\mathbb H}
\newcommand{\CC}{\mathbb C}
\newcommand{\RR}{\mathbb R}
\newcommand{\Law}{\operatorname{Law}}
\newcommand{\nuA}[4]{\nu^{#1}_{#2,#3,#4}}
\DeclareMathOperator{\Id}{Id}
\title{Research record: locality, terminal information, and planar exploration laws}
\hypersetup{pdftitle={Research record: locality, terminal information, and planar exploration laws},pdfauthor={Autonomous mathematical research}}
\author{Autonomous mathematical investigation, with independent critical checks}
\date{26 October 2026}
\begin{document}
\maketitle
\begin{abstract}
The project's three-axiom classification is false
under the usual unparameterized-curve interpretation of locality. An exact
counterexample chooses an affine SLE$_6$ metric from the boundary germ of the
terminal prime end. This germ survives every operation in the stated axioms.
The example can also be covariant under all Euclidean similarities. A
strict boundary-swallowing probability identifies an affine metric already
from one marked Jordan-domain path law, proving that no fixed affine map
normalizes the example. We also prove rigidity for deformations by one
global homeomorphism, identify precisely the extra driving-martingale
assumption that forces SLE$_6$, and record domain-continuity, target-consistency,
reversal, and near-critical obstructions. Critical percolation universality
and classification under suitably strengthened axioms remain open.
\end{abstract}
\tableofcontents

\section{Scope, conventions, and verification status}
\label{sec:scope}

This record starts from an otherwise empty research repository. The project
specification is preserved unchanged. Its primary question is whether strong
domain Markov, locality, and physical translation covariance force a
domain-indexed family of exploration curves to be one fixed affine pullback
of chordal SLE$_6$. We do not assume this conclusion or conformal covariance
of an unknown family.

Curves are oriented, unparameterized continuous paths. After an initial
exploration, the remaining domain is the component facing the terminal
prime end, with the exposed tip interpreted as an appropriate prime end.
Locality compares curves up to first encounter with a removed closed set;
their respective capacity clocks need not agree. Deletions must preserve
the source and terminal end neighborhoods and leave a simply connected
target component. This is the customary precise setting for chordal
locality. For a Jordan target and a planar closed set disjoint from that
point, preservation of its terminal neighborhood is automatic.

\paragraph{External facts used.}
We use existence, conformal covariance, strong domain Markov, locality,
and continuous Loewner generation of ordinary SLE$_6$ as established
theorems; see \cite{lawler}. We use elementary conformal mapping, boundary
extension on Jordan domains, harmonic uniqueness, and Morera's theorem.
Loewner admissibility of affine images is justified explicitly using
\cite{yuan}. A reversal statement below additionally uses \cite{ms}.
These external results are distinguished from the arguments proved here.

\paragraph{Proof checks.}
The target-germ construction was found independently in two investigations.
The boundary-race computation and the single-domain metric-identification
argument were checked independently, including the sign of the diffusion,
the exclusion of simultaneous swallowing, and the boundary-extension step.
The global-homeomorphism rigidity argument and the local boundary-twist
test were also independently checked. These checks are evidence of
correctness, not a replacement for expert human review. Nothing here proves
universality for general critical percolation.

\section{An affine metric is identifiable from one chordal law}
\label{sec:identifiability}

For a bounded Jordan domain $D$, distinct $p,q\in\partial D$, and an
orientation-preserving invertible real-linear map $A$, put
\begin{equation}
 \nuA{A}{D}{p}{q}
 = (A^{-1})_*\Law\bigl(\SLE(AD;Ap,Aq)\bigr).
 \label{eq:affine-law}
\end{equation}
Translations in an affine normalizer are immaterial. The tensor
$A^{\mathsf T}A$ is its Euclidean conformal metric, determined only up to a
positive scalar.

\subsection{A boundary-swallowing race}

\begin{lemma}[Strict boundary-race probability]
\label{lem:race}
Let $\eta$ be chordal SLE$_6$ in $\HH$ from $0$ to $\infty$. For
$x<0<y$, let $\tau_x,\tau_y$ denote their Loewner swallowing times. Then
\begin{equation}
 \mathbb P(\tau_y<\tau_x)=F\left(\frac{-x}{y-x}\right),\qquad
 F(r)=\frac{\int_0^r[u(1-u)]^{-2/3}\,du}
 {\int_0^1[u(1-u)]^{-2/3}\,du}.
 \label{eq:race}
\end{equation}
In particular $F$ is strictly increasing on $(0,1)$. The event in
\eqref{eq:race} is invariant under orientation-preserving homeomorphisms of
marked Jordan domains, when formulated as the order of separation from
the terminal prime end by the filled exploration hull.
\end{lemma}

\begin{proof}
Use $\partial_tg_t(z)=2/(g_t(z)-W_t)$ with $W_t=\sqrt6 B_t$. Until the
first swallowing time put
\[
 X_t=g_t(x)-W_t<0,\quad Y_t=g_t(y)-W_t>0,\quad
 T_t=Y_t-X_t,\quad R_t=-X_t/T_t\in(0,1).
\]
Then
\begin{align*}
 dX_t&=\frac{2}{X_t}\,dt-\sqrt6\,dB_t,&
 dY_t&=\frac{2}{Y_t}\,dt-\sqrt6\,dB_t,\\
 dT_t&=\frac{2T_t}{-X_tY_t}\,dt>0,&
 dR_t&=\frac{2(1-2R_t)}{R_t(1-R_t)T_t^2}\,dt
                 +\frac{\sqrt6}{T_t}\,dB_t.
\end{align*}
With $ds=dt/T_t^2$, the generator of $R$ is
\[
 \mathcal L=3\frac{d^2}{dr^2}
       +\frac{2(1-2r)}{r(1-r)}\frac d{dr}.
\]
The increasing solution of $\mathcal LF=0$, $F(0)=0$, $F(1)=1$ has
$F'(r)=C[r(1-r)]^{-2/3}$, giving \eqref{eq:race} by bounded optional
stopping. To justify the terminal values, $|X|/\sqrt6$ and $Y/\sqrt6$
are Bessel processes of dimension $1+4/6=5/3$ up to their first zero;
thus the first swallowing time is finite almost surely. Alternatively,
the accessible endpoints of the displayed diffusion give the same stopping
argument. There is no simultaneous swallowing, since $T_t\geq y-x>0$,
whereas both $X_t,Y_t\to0$ would force $T_t\to0$. At the first swallowing,
$R=1$ exactly when $y$ is swallowed first. ODE solutions and the continuous
driver stay finite up to that first time.

Swallowing means that the marked boundary end has been cut off from the
target by the explored filled hull (including the boundary swallowed at
first contact). This definition uses only the domain, the oriented path,
and its order of traversal. A homeomorphism carries filled target
components to filled target components. Capacity parametrization therefore
has no role in the comparison event.
\end{proof}

\subsection{Boundary rigidity and affine uniqueness}

\begin{lemma}[A single SLE law identifies its conformal boundary map]
\label{lem:boundary-rigidity}
Let $h:\overline D\to\overline{D'}$ be an orientation-preserving
homeomorphism between closed bounded Jordan domains, with
$p'=h(p)$, $q'=h(q)$. If $h_*\SLE(D;p,q)=\SLE(D';p',q')$, then
$h|_{\partial D}$ agrees with the boundary extension of a conformal
map $f:D\to D'$.
\end{lemma}

\begin{proof}
Choose conformal maps $\phi:D\to\HH$ and $\phi':D'\to\HH$ sending
their marked endpoints to $0,\infty$. Their Jordan boundary extensions
give an increasing homeomorphism
\[
 H=\phi'\circ h\circ\phi^{-1}:\overline\RR\to\overline\RR
\]
fixing $0,\infty$. For every $x<0<y$, equality of the corresponding
boundary-race probabilities and strictness in Lemma~\ref{lem:race} give
\[
 \frac{-H(x)}{H(y)-H(x)}=\frac{-x}{y-x},\quad\text{hence}\quad
 \frac{H(x)}{H(y)}=\frac{x}{y}.
\]
Fixing one $x<0$ and one $y>0$ shows $H(s)=cs$ for all real $s$, with
$c>0$. Consequently $f=(\phi')^{-1}\circ(c\phi)$ has exactly the
required boundary values. Equality of interior maps is not asserted for
an arbitrary $h$.
\end{proof}

\begin{theorem}[Single-domain affine identifiability]
\label{thm:affine-identifiability}
For even one bounded Jordan $D$ and one pair of distinct boundary points
$p,q$, the equality $\nuA{A}{D}{p}{q}=\nuA{B}{D}{p}{q}$ holds if and
only if $BA^{-1}$ is a Euclidean similarity. Equivalently,
\[
 B^{\mathsf T}B=c A^{\mathsf T}A\qquad(c>0).
\]
The statement also allows orientation-reversing normalizers by composing
one map with a reflection.
\end{theorem}

\begin{proof}
The equality of laws means $M=BA^{-1}$ carries the SLE law in $AD$ to
the SLE law in $BD$, at the corresponding endpoints. By
Lemma~\ref{lem:boundary-rigidity}, a conformal map $f:AD\to BD$ has
boundary values equal to $M$. Both real coordinate functions of the
real-linear map $M$ are harmonic. The real and imaginary parts of $f$
are also harmonic, continuous on the closure, with the same boundary
values. The maximum principle on the bounded domain gives $f=M$
throughout $AD$. Thus $M$ is complex-linear and invertible, hence a
similarity. The converse is ordinary SLE similarity covariance.
Reflection covariance reduces the remaining orientation case to this one.
\end{proof}

\section{An exact counterexample using the terminal boundary germ}
\label{sec:target-germ}

\subsection{The conserved information and the admissible domain class}

Let $G(D,q)$ be the embedded germ of $D$ at its terminal prime end.
Concretely, target neighborhoods may be taken as
\[
 f^{-1}\bigl(\HH\setminus\overline{B(0,R)}\bigr),\qquad R\to\infty,
\]
for any conformal $f:D\to\HH$ sending $q$ to infinity. The embedded
germ records the planar domain in a full neighborhood of that end;
one may shrink neighborhoods without changing it. A full neighborhood
contains a neighborhood in the prime-end compactification, rather than
merely an arbitrarily narrow approach sector. Germ equivalence means
agreement on such neighborhoods with the same inherited marked end.

At a Jordan boundary point this is the ordinary local boundary germ.
At a slit tip one must retain which prime end is intended. A germ is
called $C^1$-regular here when it has a singleton endpoint and, in a full
end neighborhood, a single $C^1$ boundary arc with a nonzero tangent and
the domain locally on one side. The tangent line $\ell(G)$ is unoriented.
All other germs receive the default label below. Regularity and tangent
are properties of the germ itself.

Finite chordal Loewner hulls in a target-normalized half-plane chart are
bounded. They therefore miss some target neighborhood. At every finite
stopping time prior to reaching the target, the remaining marked domain
$D_t$ satisfies
\begin{equation}
 G(D_t,q)=G(D,q).
 \label{eq:germ-persistent}
\end{equation}
Likewise locality deletions away from the terminal end preserve this germ.
Thus the following proof already works on domains with regular targets,
together with their slit domains, a class closed under all operations
needed in the argument. Assigning defaults to nonregular end germs extends
it to arbitrary prime ends under the stated end-neighborhood convention.
Locality for sets accumulating at a nonsingleton impression must always be
formulated in prime-end terms; it is not defined by pretending the end is
a single physical point.

\begin{proposition}[Conserved-label construction]
\label{prop:germ-construction}
Assign to each embedded terminal germ an orientation-preserving invertible
real-linear map $A(G)$, measurably and in a translation-covariant way:
$A(G(D+v,q+v))=A(G(D,q))$. Define
\begin{equation}
 P_{D,p,q}=(A(G(D,q))^{-1})_*
       \SLE\bigl(A(G(D,q))D;A(G(D,q))p,A(G(D,q))q\bigr).
 \label{eq:germ-ensemble}
\end{equation}
Then this family has strong domain Markov, locality, and translation
covariance as unparameterized paths.
\end{proposition}

\begin{proof}
Write $A=A(G(D,q))$. By \eqref{eq:germ-persistent} the same deterministic
$A$ is prescribed after any finite initial exploration. Apply strong
domain Markov to SLE$_6$ in $AD$, then pull the tail back by $A^{-1}$.
This gives exactly \eqref{eq:germ-ensemble} in the remaining target
component, at its exposed tip and target. Stopping times are carried
to stopping times because the invertible map preserves the explored
filtration; the capacity clock is an adapted increasing time change.
For deletion of $K$ away from both ends, the two domains use the same
$A$. SLE$_6$ locality in $AD$ and $AD\setminus AK$, pulled back by
$A^{-1}$, couples their paths identically until first encounter with $K$.
Finally the displayed covariance of the label and translation covariance
of SLE$_6$ prove physical translation covariance. No conditioning on a
hidden random parameter occurs: the label is a deterministic conserved
function of the current marked domain.
\end{proof}

\subsection{A concrete example with full similarity covariance}

Fix $a>0$, $a\ne1$, and set
\begin{equation}
 A(G)=\begin{cases}
 \Pi_{\ell(G)}+a\Pi_{\ell(G)^\perp},&G\text{ is }C^1\text{-regular},\\
 I,&\text{otherwise}.
 \end{cases}
 \label{eq:tangent-label}
\end{equation}
Here $\Pi$ is orthogonal projection. On smooth-target domains this is a
continuous function of the tangent in the topology controlling that
tangent; on arbitrary germs it is a deterministic local geometric rule.
If desired, the class can first be restricted to piecewise smooth
domains and their finite slit descendants, avoiding any measurable-space
conventions for the class of all planar domains.

For a similarity $S(x)=cRx+v$, $c>0$, $R$ orthogonal, one has
\[
 A(G(SD,Sq))=R A(G(D,q))R^{-1}.
\]
Combining this with SLE$_6$ similarity and reflection covariance proves
$P_{SD,Sp,Sq}=S_*P_{D,p,q}$. Thus adding scale invariance, reflection
covariance, or covariance under every rotation does not exclude this
construction.

\begin{theorem}[The proposed fixed-affine classification is false]
\label{thm:counterexample}
The ensemble \eqref{eq:germ-ensemble}--\eqref{eq:tangent-label} satisfies
the three proposed axioms, and full Euclidean similarity covariance, but
there is no one fixed affine transformation sending all its laws to
ordinary chordal SLE$_6$. This failure already occurs for smooth targets
on one disk.
\end{theorem}

\begin{proof}
Take $D$ to be the unit disk and fix $p=-1$. At $q_h=i$ the tangent is
horizontal, while at $q_v=1$ it is vertical. The respective matrices are
\[
 A_h=\begin{pmatrix}1&0\\0&a\end{pmatrix},\qquad
 A_v=\begin{pmatrix}a&0\\0&1\end{pmatrix}.
\]
A fixed affine normalizer with linear part $B$ would imply
$\nuA{B}{D}{p}{q_h}=\nuA{A_h}{D}{p}{q_h}$ and
$\nuA{B}{D}{p}{q_v}=\nuA{A_v}{D}{p}{q_v}$.
Theorem~\ref{thm:affine-identifiability} makes both $BA_h^{-1}$ and
$BA_v^{-1}$ similarities. Hence $A_vA_h^{-1}=\operatorname{diag}(a,a^{-1})$
would be a similarity, forcing $a^2=1$, contrary to $a>0$, $a\ne1$.
\end{proof}

\subsection{Why these really are continuous-driving Loewner curves}
\label{sec:loewner-admissible}

Continuity and the absence of transverse crossings alone would not suffice
to settle this point. Instead use the deterministic characterization in
\cite[Theorem 1.1 and Remark 1.4]{yuan}: a continuous half-plane path with
strictly increasing filled hulls is a continuous-driving Loewner trace if
and only if it is a locally uniform limit of simple interior paths. The
condition is independent of parametrization.

An ordinary SLE$_6$ trace has this property. Conjugating an invertible
affine map by the conformal charts of its source and image domains gives
a quasiconformal self-map of $\HH$. It extends to a homeomorphism of the
closed half-plane with its point at infinity, fixing the two marked ends
after normalization. It carries simple interior approximants to simple
interior approximants; uniform convergence survives on compact time
intervals because the extended map is uniformly continuous on relevant
compact sets. Filled hulls and their strict growth are topologically
preserved. Yuan's theorem gives a continuous driving function after
reparametrizing by capacity. This also treats nonsmooth slit domains in
the Markov property without assuming Euclidean continuity of their
uniformizing maps on every boundary point. The counterexample therefore
applies to the driving-function formulation as well as to curves.

% Technical completion of affine-image Loewner admissibility.
\subsection{Quasiconformal images have infinite capacity lifetime}
\label{sec:qc-capacity-lifetime}

Yuan's characterization gives a continuous Loewner driver after recoding
an affine image by its capacity, but this alone does not assert that the
capacity lifetime is infinite.  A continuous curve tending to infinity need
not have infinite half-plane capacity.  The following argument proves the
required assertion for the quasiconformal changes arising here.

\begin{lemma}
Let $F:\mathbb H\to\mathbb H$ be a fixed $K$-quasiconformal
self-homeomorphism whose boundary extension fixes $0$ and $\infty$.
If $\eta$ is ordinary chordal SLE$_6$ from $0$ to $\infty$, then
$F\eta$ has unbounded imaginary part almost surely.  Consequently its
half-plane-capacity lifetime is infinite.
\end{lemma}

\begin{proof}
Put $C=\overline{B(i,1/4)}\subset\mathbb H$.  The boundary map of $F$
is increasing, so $F(R)>0$ for $R>0$ and $F(R)\to+\infty$ as
$R\to\infty$.  The normalized maps
\[
 F_R(z)=\frac{F(Rz)}{F(R)}
\]
are $K$-quasiconformal half-plane maps fixing $0,1,\infty$.
Reflecting them across the real line gives normalized quasiconformal sphere
maps.  Compactness of this normalized family, including its inverses,
implies
\[
 \inf_{R>0}\inf_{z\in C}\operatorname{Im}F_R(z)=c_K>0.
\]
Indeed, a sequence violating positivity would have a subsequence converging
uniformly on the sphere to a quasiconformal homeomorphism preserving the
half-plane, and a subsequence of points in $C$ would then be mapped to its
real boundary, an impossibility.  Hence
\[
 \inf_{z\in RC}\operatorname{Im}F(z)\geq c_K F(R)
       \longrightarrow\infty.
 \tag{QC1}
\]

The event $E=\{\eta\cap C\ne\varnothing\}$ has positive probability.
One elementary justification avoids invoking a support theorem.  In the
normalization $\partial_tg_t=2/(g_t-W_t)$, the Brownian event
\[
 \sup_{0\leq t\leq1}|W_t|<\varepsilon,
       \qquad 0<\varepsilon<1/4,
\]
has positive probability.  On this event the Loewner hull stays within the
vertical strip $|\operatorname{Re}z|\leq\varepsilon$ up to time $1$.
For example, a point initially to the right of that strip has increasing
real part under the Loewner ODE and retains a positive real gap from the
driver, so it cannot be swallowed in finite time; the other side is
identical.  If the entire trace up to time $1$ had height at most $1$, its
filled hull would lie in the half-disk of radius
$\sqrt{1+\varepsilon^2}$.  By capacity monotonicity its capacity would be
at most $1+\varepsilon^2<2$, whereas its actual capacity is $2$.
Here a half-disk of radius $R$ has hydrodynamic map $z+R^2/z$ and
capacity $R^2$.  Thus the continuous trace crosses height $1$ in the strip,
and that crossing lies in $C$.

Fix $r>1$.  Brownian scaling, and therefore scalar dilation of ordinary
SLE$_6$, is mixing.  Indeed $Z_s=e^{-s/2}B_{e^s}$, $s\in\mathbb R$,
is a stationary Gaussian process with covariance $e^{-|s-u|/2}$; scalar
Brownian scaling becomes a log-time shift.  Decay of these covariances
gives mixing first for Gaussian cylinder functions and then for the whole
path sigma-field.  Its inverse dilation is ergodic as well.  Birkhoff's
theorem applied to $1_E$ gives
\[
 \eta\cap r^nC\ne\varnothing
       \qquad\text{for infinitely many positive integers }n,
       \quad\text{almost surely}.
\]
In fact these visits have limiting positive frequency, since
$1_E(S_{r^{-n}}\eta)=1_{\{\eta\cap r^nC\ne\varnothing\}}$.
Equation (QC1) then shows that $F\eta$ has unbounded height.

Every finite initial image trace is compact away from infinity and has
finite capacity.  The standard Loewner height estimate
\[
 \max\{\operatorname{Im}z:z\in K\}^2
       \leq2\operatorname{hcap}(K)
\]
for a half-plane hull, or equivalently
$\max\operatorname{Im}K_u\leq2\sqrt u$ when
$\operatorname{hcap}(K_u)=2u$, now forces the capacities of these image
prefixes to tend to infinity.  Thus the capacity-recoded image has lifetime
$[0,\infty)$.
\end{proof}

For an affine image in a bounded marked domain, conjugating the affine map
by its source and image conformal charts gives exactly such a fixed
quasiconformal $F$ fixing the marked ends.  This lemma therefore completes
the infinite-time driving-function assertion for the target-germ
counterexample.  It uses properties of SLE$_6$, rather than an invalid
general assertion about arbitrary curves tending to infinity.

\section{A positive classification within the global-deformation class}
\label{sec:global-rigidity}

The preceding example does not arise from one deterministic global change
of spatial coordinates. Within that narrower class the proposed affine
conclusion does follow, and no differentiability of the change of
coordinates is required.

\begin{theorem}[Full-family homeomorphism symmetries are conformal]
\label{thm:homeo-symmetry}
Let $h:\CC\to\CC$ be an orientation-preserving homeomorphism such that,
for every bounded Jordan domain and marked boundary pair,
\[
 h_*\SLE(D;p,q)=\SLE(hD;hp,hq).
\]
Then $h$ is an entire conformal automorphism, hence $h(z)=\alpha z+\beta$
for complex $\alpha\ne0$. For an orientation-reversing homeomorphism,
the corresponding conclusion is anticonformality.
\end{theorem}

\begin{proof}
Take any nondegenerate closed Euclidean triangle $T$ and distinct
$p,q\in\partial T$. Lemma~\ref{lem:boundary-rigidity} gives a conformal
$f_T:T^\circ\to h(T^\circ)$ whose continuous boundary extension equals
$h$ on $\partial T$. Cauchy's theorem, applied on interior triangles
and then passed to the boundary by continuity, yields
\[
 \int_{\partial T}h(z)\,dz=\int_{\partial T}f_T(z)\,dz=0.
\]
Since $h$ is continuous and this holds for every triangle, Morera's
theorem makes $h$ entire holomorphic. Its injectivity makes it conformal,
and an entire conformal automorphism of the plane is complex affine.
Composition with a reflection proves the other orientation case.
\end{proof}

\begin{theorem}[Stationary global deformations are real affine]
\label{thm:stationary-deformation}
Let $h:\CC\to\CC$ be a homeomorphism and define, on all bounded Jordan
domains,
\[
 P^h_{D,p,q}=(h^{-1})_*\SLE(hD;hp,hq).
\]
If this ensemble is covariant under every physical translation, then
$h$ is a real-affine homeomorphism.
\end{theorem}

\begin{proof}
Reflection covariance lets us assume $h$ preserves orientation. For every
$v\in\CC$, physical translation covariance says that
\[
 \psi_v=h\circ T_v\circ h^{-1},\qquad T_v(z)=z+v,
\]
preserves the full chordal SLE$_6$ family. All bounded Jordan domains are
available, since inverse images under a plane homeomorphism are bounded
Jordan domains. Theorem~\ref{thm:homeo-symmetry} gives
$\psi_v(z)=\alpha_v z+\beta_v$. For $v\ne0$, the conjugate of a nonzero
translation has no fixed point. A complex affine map with
$\alpha_v\ne1$ has a fixed point, so $\alpha_v=1$ for all $v$.
It follows that $h(x+v)=h(x)+\beta_v$, and putting $x=0$ gives
\[
 h(x+v)-h(0)=(h(x)-h(0))+(h(v)-h(0)).
\]
The continuous additive map $h-h(0)$ is real-linear. It is invertible
because $h$ is a homeomorphism.
\end{proof}

\paragraph{What this does not prove.}
The original axioms do not give a representation by one global $h$.
The counterexample shows that replacing a domain-dependent geometry by
one fixed spatial geometry is a substantive missing step, not an
innocent coordinate choice.

\section{Boundary stability and other ways to remove terminal dependence}
\label{sec:boundary-stability}

\subsection{An explicit failure of continuity under small target perturbations}

The tangent rule is continuous when target tangents converge, but not under
ordinary uniform approximation of the domains and their boundaries.
This defect persists even when every domain is smooth and the approximating
domains are uniform quasidisks.

\begin{proposition}[Local rotation of a terminal germ]
\label{prop:boundary-discontinuity}
For the ensemble of Theorem~\ref{thm:counterexample}, there are smooth
Jordan $D_n\to D$, with uniformly converging boundary parametrizations
and fixed distinct marked points $p,q$, for which $P_{D_n,p,q}$ does
not converge to $P_{D,p,q}$ in the uniform unparameterized-curve topology.
\end{proposition}

\begin{proof}
Let $D$ be the unit disk, $q=i$, $p=-1$. Choose a smooth radial cutoff
$\chi$ equal to $1$ near $0$ and $0$ outside the unit disk, and define
the orientation-preserving smooth diffeomorphism
\[
 F(re^{i\theta})=r\exp\bigl(i\theta+i(\pi/2)\chi(r)\bigr).
\]
It is a rotation by $\pi/2$ near $0$ and the identity outside the disk.
Its quasiconformal distortion is finite, since $r\chi'(r)$ is bounded.
Put $F_\epsilon(z)=q+\epsilon F((z-q)/\epsilon)$ and
$D_\epsilon=F_\epsilon(D)$. Then $F_\epsilon(q)=q$, and $F_\epsilon(p)=p$
for small $\epsilon$, while $\sup_z|F_\epsilon(z)-z|\leq2\epsilon$.
All $F_\epsilon$ have one common quasiconformal bound. At $q$ the target
tangent has rotated by $\pi/2$, so the matrix used for every
$D_\epsilon$ is $A_v$, whereas the matrix for $D$ is $A_h$.

The domains $A_vD_\epsilon$ are uniform quasidisks converging to $A_vD$,
with the marked boundary points converging as well. Kernel convergence
and uniform boundary equicontinuity of normalized Riemann maps on
uniform quasidisks give uniform convergence of these maps on the closed
disk. Coupling their chordal SLE$_6$ paths with one common disk SLE path
therefore gives convergence of the entire unparameterized curves. Thus
\[
 P_{D_\epsilon,p,q}\ \Longrightarrow\ \nuA{A_v}{D}{p}{q}.
\]
By Theorem~\ref{thm:affine-identifiability}, this limit differs from
$P_{D,p,q}=\nuA{A_h}{D}{p}{q}$ because
$A_vA_h^{-1}=\operatorname{diag}(a,a^{-1})$ is not a similarity.
\end{proof}

This proves that an appropriate boundary-stability requirement excludes
the example. It does not prove that boundary stability plus the three
original axioms suffices for classification.

\begin{corollary}[Boundary stability classifies tangent-dependent affine laws]
\label{cor:tangent-stability}
Suppose, on smooth bounded Jordan domains with smooth targets, an ensemble
has the form $P_{D,p,q}=\nuA{A_{\ell(D,q)}}{D}{p}{q}$, where $A_\ell$
is an arbitrary orientation-preserving invertible linear map depending
only on the unoriented target tangent line. If these laws are continuous
under the uniform smooth-boundary approximations used in
Proposition~\ref{prop:boundary-discontinuity}, then one fixed affine
map normalizes all of them. If the ensemble is also covariant under
all rotations, that fixed map is a similarity.
\end{corollary}
\begin{proof}
Given any two tangent lines $\ell_1,\ell_2$, take a disk having tangent
$\ell_1$ at $q$ and apply the preceding local twist with the rotation
carrying $\ell_1$ to $\ell_2$. Every approximating domain uses the
constant matrix $A_{\ell_2}$. SLE continuity in that fixed metric shows
that the limiting law is $\nuA{A_{\ell_2}}{D}{p}{q}$, whereas the assumed
boundary stability makes it $\nuA{A_{\ell_1}}{D}{p}{q}$. Single-pair
identifiability implies $A_{\ell_2}A_{\ell_1}^{-1}$ is a similarity.
Thus all $A_\ell^{\mathsf T}A_\ell$ are proportional, which is exactly
one fixed affine family. For rotation covariance its fixed tensor $Q$
satisfies $R^{\mathsf T}QR=Q$ for every rotation $R$, by the same metric
uniqueness, hence $Q$ is scalar.
\end{proof}

\subsection{Reversibility eliminates this mechanism inside the affine-germ ansatz}

Here is a rigorous limitation of the constructed counterexamples, useful for
formulating a more faithful percolation classification.  Suppose
\[
 P_{D,p,q}=A(D,q)^{-1}_*\mathrm{SLE}_6(AD;Ap,Aq)
\]
on bounded Jordan domains, and suppose this family is reversible:
reversing $P_{D,p,q}$ gives $P_{D,q,p}$.  Then, for all $p,q$,
\[
 A(D,p)^T A(D,p)=c(D,p,q)A(D,q)^T A(D,q),\qquad c(D,p,q)>0.
\]
Use here the external reversibility theorem of J. Miller and S. Sheffield,
\emph{Imaginary geometry III: reversibility of SLE$_\kappa$ for
$\kappa\in(4,8)$}, Annals of Mathematics 184 (2016), 455--486,
\url{https://annals.math.princeton.edu/2016/184-2/p03}.
Indeed, ordinary SLE$_6$ reversibility says that reversing the first curve
uses the same matrix $A(D,q)$, now with endpoints $q,p$.  By assumed
reversibility this is equal to the law constructed with $A(D,p)$.
Single-pair affine identifiability therefore says that
$A(D,p)A(D,q)^{-1}$ is a similarity, which is equivalent to the displayed
matrix relation.  For a rectangle with one endpoint in a horizontal side
and the other in a vertical side, the proposed matrices
$\operatorname{diag}(1,a)$ and $\operatorname{diag}(a,1)$ violate this relation
when $a\ne1$.  Thus the construction explicitly fails reversibility.

This result only eliminates endpoint-dependent anisotropy \emph{within this
ansatz}; it does not prove the desired classification after reversibility
is added.  Near-critical percolation suggests a separate criticality
condition remains necessary.

\section{Further affine rigidity and target consistency}
The following argument concerns ordinary chordal SLE$_6$; it does not
assume conformal invariance of an unknown exploration family.
For SLE$_6$ in $\mathbb H$ from $0$ to $\infty$, the left-passage event $L(z)$
is defined by the eventual side of the oriented trace containing $z$.
For a nonsimple trace this is also its eventual Loewner swallowing side.
It is invariant under orientation-preserving homeomorphisms of the marked
half-plane. A fixed interior point belongs to the trace with probability zero.

\begin{lemma}[Strict left-passage function]
\label{lem:am-left}
For $z=x+iy$, $y>0$,
\[
 \mathbb P[L(z)]=F(x/y),\qquad
 F(u)=\frac{\int_{-\infty}^u(1+v^2)^{-2/3}\,dv}
 {\int_{-\infty}^{\infty}(1+v^2)^{-2/3}\,dv}.
\]
In particular, $F$ is strictly increasing.
\end{lemma}
\begin{proof}
This is the standard SLE left-passage formula; its short diffusion derivation
explains why strict monotonicity holds even for the nonsimple SLE$_6$ trace.
Write $g_t(z)-\sqrt\kappa B_t=X_t+iY_t$. Then
\[
 dX_t=\frac{2X_t}{X_t^2+Y_t^2}\,dt-\sqrt\kappa\,dB_t,
 \qquad dY_t=-\frac{2Y_t}{X_t^2+Y_t^2}\,dt.
\]
For $U=X/Y$ and $ds=dt/Y_t^2$, It\^o's formula gives
\[
 dU_s=\frac{4U_s}{1+U_s^2}\,ds-\sqrt\kappa\,d\widetilde B_s.
\]
The increasing scale function with limits $0,1$ solves
\[
 \frac\kappa2 F''+\frac{4u}{1+u^2}F'=0,
 \qquad F'(u)=C_\kappa(1+u^2)^{-4/\kappa}.
\]
The normalizing integral is finite for $\kappa<8$. The clock runs to
infinity up to the point's maximal Loewner lifetime: otherwise
$d\log Y_s=-2(1+U_s^2)^{-1}ds$ keeps $Y$ positive and the diffusion extends.
The bounded martingale $F(U_s)$ converges. A limit in $(0,1)$ would force
$U_s$ to converge to a finite value and its martingale quadratic variation
to be infinite, a contradiction. Thus $U_s\to-\infty$ or $+\infty$.
Optional stopping and the usual topological side identification for
chordal Loewner chains give the formula at $\kappa=6$.
\end{proof}

\begin{proposition}[Linear symmetries of half-plane SLE$_6$]
\label{prop:am-linear}
If an invertible real-linear map $M$ carries SLE$_6(\mathbb H,0,\infty)$
to SLE$_6(M\mathbb H,0,\infty)$, then $M$ is a Euclidean similarity,
possibly reversing orientation.
\end{proposition}
\begin{proof}
Reflection invariance reduces to $\det M>0$. Rotate $M\mathbb H$ back to
$\mathbb H$. The resulting map is
$T(x,y)=(ax+by,cy)$ with $a,c>0$. By topological preservation of the side
event and Lemma~\ref{lem:am-left},
\[
 F(u)=F((au+b)/c)\qquad(u\in\mathbb R).
\]
Strict monotonicity implies $a=c$, $b=0$. Undo the rotation.
\end{proof}

\subsection{Target consistency is a substantive additional condition}

\paragraph{Additional axiom; not a consequence proved here.}
For two targets $q_1,q_2$ on a boundary arc $J$ whose closure avoids $p$,
require explorations towards them to couple identically until first
hitting $J$ (or use the usual stopping formulation before target separation).
A common percolation coloring gives this: the two boundary conditions
disagree only on $J$, so the exploration algorithms see identical colors
until reaching it. The stated locality axiom fixes $q$ and does not assert
this property.

\begin{proposition}[Target consistency removes variable affine target metrics]
\label{prop:am-targetconsistency}
Suppose $D$ is a bounded Jordan domain analytic near $p$, and
\[
 \mathcal P_{D,p,q}=(A_q^{-1})_*\operatorname{SLE}_6(A_qD,A_qp,A_qq).
\]
If target consistency holds for $q_1,q_2$ on an arc avoiding $p$, then
$A_{q_2}A_{q_1}^{-1}$ is a similarity. In particular, their metric tensors
$A_q^{\mathsf T}A_q$ are proportional.
\end{proposition}
\begin{proof}
Translate $p$ to zero and dilate by $r\to\infty$. The domains converge
locally to the tangent half-plane $H_p$, and targets tend to infinity.
For each fixed $q$, conformal kernel convergence at the analytic start
and the usual initial-segment construction of SLE give limiting law
\[
 (A_q^{-1})_*\operatorname{SLE}_6(A_qH_p,0,\infty).
\]
Convergence can be tested on paths stopped on crosscuts strictly inside
fixed half-disks, followed by exhaustion. The dilated disagreement arc
escapes every fixed neighborhood, so target consistency says that the
two limiting half-plane laws coincide. Apply $A_{q_1}$, rotate the image
half-plane, and use Proposition~\ref{prop:am-linear}.
\end{proof}

This eliminates the specific target-tangent obstruction, not every
possible counterexample. The proof uses initial-segment convergence for
ordinary SLE$_6$, not for an unknown general ensemble.

\subsection{A compatibility graph for all target-germ affine decorations}

The previous target-consistency proposition also gives a useful classification
within the entire class of affine decorations, without a claim about arbitrary
exploration ensembles. Define a graph whose vertices are admissible embedded
target germs, modulo translations, and join $G_1,G_2$ when they can both occur
on one bounded Jordan domain with an analytic starting point $p$ elsewhere.
The targets must lie on a boundary arc whose closure avoids $p$.
Assign to each vertex the projective positive-definite tensor
$[Q(G)]=[A(G)^{\mathsf T}A(G)]$.

\begin{corollary}[Target-germ metrics are constant on compatibility components]
\label{cor:am-germgraph}
For a target-germ affine SLE$_6$ decoration satisfying target consistency,
$[Q(G)]$ is constant on every connected component of this graph.
If the graph is connected, the decoration is one fixed-affine SLE$_6$ family.
\end{corollary}
\begin{proof}
Proposition~\ref{prop:am-targetconsistency} makes the tensors proportional
across every edge. Propagate along finite paths.
If a component has representative matrix $A_0$, proportionality implies
$A(G)=c_G R_G A_0$ with $c_G>0$ and $R_G$ orthogonal.
Similarity covariance of SLE$_6$ then gives
\[
 A(G)^{-1}_*\operatorname{SLE}_6(A(G)D,A(G)p,A(G)q)
 = (A_0^{-1})_*\operatorname{SLE}_6(A_0D,A_0p,A_0q),
\]
independently of $G$.
\end{proof}

For ordinary smooth local boundary germs, connectivity is suggested by a
simple planar surgery: place two small representatives far apart, retain
their marked local boundary arcs, join the remaining endpoints with disjoint
arcs to form one Jordan boundary, and place an analytic start on a joining
arc. A fully general surgery statement for nonlocally connected prime-end
impressions is not asserted here; the graph formulation keeps this issue
explicit. For the straight-line tangent germs in the counterexample, joint
realization on an ordinary rounded rectangle already connects horizontal
and vertical targets and excludes their distinct tensors.
\subsection{Finite rotations and a fixed anisotropy}

\begin{proposition}
\label{prop:am-rotation}
Suppose an ensemble has one fixed affine SLE$_6$ normalization with linear
part $A$, and is covariant under rotation $R_\theta$, where
$\theta\notin\pi\mathbb Z$. Then $A^{\mathsf T}A$ is scalar, so $A$ is a
similarity. A half-turn alone does not suffice.
\end{proposition}
\begin{proof}
The map $AR_\theta A^{-1}$ preserves the SLE$_6$ family and hence is a
similarity. Its determinant is $1$, so it is orthogonal. For
$Q=A^{\mathsf T}A$ this says $R_\theta^{\mathsf T}QR_\theta=Q$.
A symmetric real $2\times2$ matrix commuting with a rotation other than
$0$ or $\pi$ is scalar. For $R_\pi=-I$, every $Q$ satisfies the relation.
\end{proof}

The fixed-normalizer hypothesis is essential: the target-tangent
construction can be covariant under every rotation and still have no
fixed normalizer. Rotation symmetry by itself does not repair the conjecture.


% Research fragment; intended for inclusion in progress.tex (amsmath/amssymb/amsthm).
\section{What locality actually says about a semimartingale driving process}
\label{sec:loewner-locality}

\paragraph{Status.}
The coordinate-change identities below are proved, under explicitly stated
continuous-semimartingale assumptions.  The conditional martingale theorem
does not require absolute continuity of quadratic variation.  They give a conditional SLE
classification.  The original three axioms do not currently imply these
semimartingale or local-martingale assumptions, and no such implication is
claimed.  In particular the domain Markov property does not make a driving
process Markov in its single real coordinate.

\subsection{Deletion of a hull and the drift anomaly}

Use chordal Loewner normalization
\[
 \partial_t g_t(z)=\frac{2}{g_t(z)-W_t},\qquad
 g_t(z)=z+\frac{2t}{z}+O(z^{-2}).
\]
Suppose on an interval before collision with a fixed hull that
\[
 dW_t=dM_t+b_t\,dt,\qquad d\langle M\rangle_t=v_t\,dt,
 \tag{L1}
\]
where $M$ is a continuous local martingale and $v_t\geq0$.
Let $A$ be a bounded half-plane hull away from the initial tip, and let
$h_0:\mathbb H\setminus A\to\mathbb H$ have hydrodynamic normalization at
infinity.  Grow the same trace in the deleted domain, mapping by $h_0$ first,
and denote its Loewner maps by $\widehat g_u$.  Up to the first encounter with
$A$, write
\[
 h_t=\widehat g_{u(t)}\circ h_0\circ g_t^{-1},\qquad
 \widehat W_{u(t)}=h_t(W_t),\qquad
 \dot u(t)=h_t'(W_t)^2.
 \tag{L2}
\]
The map $h_t$ is analytic across the real line near $W_t$, by Schwarz
reflection, and $h_t'(W_t)>0$.  These formulas are local near the driving
point; no differentiability of the trace is used.

Differentiating the composition in (L2) gives
\[
 \partial_t h_t(z)=
 \frac{2h_t'(W_t)^2}{h_t(z)-h_t(W_t)}
 -\frac{2h_t'(z)}{z-W_t}.
 \tag{L3}
\]
For $z=W_t+\epsilon$ the first term is
$2h_t'(W_t)/\epsilon-h_t''(W_t)+O(\epsilon)$,
while the second term is
$2h_t'(W_t)/\epsilon+2h_t''(W_t)+O(\epsilon)$.
Consequently
\[
 (\partial_t h_t)(W_t)=-3h_t''(W_t).
 \tag{L4}
\]
The random field $h_t(z)$ has finite variation in time locally before
collision, so ordinary It\^o's formula (equivalently the finite-variation
case of It\^o--Wentzell) applies to $h_t(W_t)$.  Thus
\[
 d\widehat W_{u(t)}=
 h_t'(W_t)\,dM_t+
 \left[h_t'(W_t)b_t+
        \left(\frac{v_t}{2}-3\right)h_t''(W_t)\right]dt.
 \tag{L5}
\]
After the capacity time change, its characteristics are therefore
\[
 \boxed{\quad
 \widehat v_{u(t)}=v_t,\qquad
 \widehat b_{u(t)}=
   \frac{b_t}{h_t'(W_t)}+
   \left(\frac{v_t}{2}-3\right)
       \frac{h_t''(W_t)}{h_t'(W_t)^2}.
 \quad}
 \tag{L6}
\]
The coefficient $-3$, and hence the exceptional variance $6$, follow from
the normalization $\partial_tg_t=2/(g_t-W_t)$.  A different capacity
normalization changes both constants consistently.

\paragraph{Logical content.}
Locality identifies the law of the transformed driving process with the
law prescribed for the deleted physical domain in its own conformal chart.
It does \emph{not} identify that law with the original driving law in its
original chart.  In particular the two drifts in (L6) need not vanish,
or be equal.  A geometry-dependent drift can transform with the second
term in (L6); deleting that term without an additional hypothesis is exactly
where an invalid proof of $v=6$ would enter.  Physical translation invariance
acts by translating the physical domain and trace.  It supplies no invariance
under the nonlinear map $h_t$ and cannot by itself justify identifying these
two characteristics.

\subsection{A precise conditional classification by driving martingales}

\paragraph{Conditional theorem (proved).}
\leavevmode\par\noindent
Suppose an ensemble of continuous Loewner-generating chordal traces satisfies
strong domain Markov and locality, with the conventional endpoint-preserving
hull deletions.  Suppose also that, for every deterministic marked domain,
its capacity driving process in a chordal conformal chart is a continuous
local martingale in its natural filtration.  Assume the usual measurability
and restart consistency of this family, so that the strong domain Markov
property applies in the chart at a random slit domain.  Then
\[
 W_t-W_0=\sqrt6\,B_t,
\]
and the ensemble is ordinary chordal $\mathrm{SLE}_6$.
No absolute continuity or differentiability assumption on quadratic
variation is needed.  The chart premise is independent of which chordal
chart is selected, since two such charts differ by $z\mapsto az+b$, $a>0$,
and the corresponding deterministic capacity time change preserves local
martingales.

\begin{proof}
Fix an initial marked domain and map it to $(\mathbb H,0,\infty)$.
Delete the inverse image of a vertical slit at $x\ne0$, of height $2\sqrt s$
for $s>0$.  The hydrodynamic map of this slit is
\[
 h_0(z)=x+\sqrt{(z-x)^2+4s},
 \tag{L7}
\]
with the branch asymptotic to $z-x$ at infinity.  At zero,
\[
 h_0'(0)=\frac{|x|}{\sqrt{x^2+4s}},\qquad
 h_0''(0)=-\frac{4s\,\operatorname{sgn}(x)}{(x^2+4s)^{3/2}}\ne0.
 \tag{L8}
\]
Let $C_t=\langle W\rangle_t$, initially with no regularity beyond the
continuity and monotonicity of quadratic variation.  The general
It\^o version of (L5) is
\[
 d\widehat W_{u(t)}
   =h_t'(W_t)\,dW_t
       +\frac12h_t''(W_t)\,dC_t-3h_t''(W_t)\,dt.
 \tag{L8a}
\]
By locality and the hypothesis on the deleted domain,
$\widehat W_u$ is a continuous local martingale; its time change $\widehat W_{u(t)}$ is a local martingale in the same trace filtration. Indeed, the inverse clock $t(u)$ is the original-domain capacity of the initial trace reconstructed from the deleted-domain curve. It is an adapted, continuous increasing functional of that trace, so $u(t)$ is a stopping time in the deleted-driver filtration. The usual stopping-time change theorem, after localization, therefore applies.
That filtration is also the one generated by $W_t$, since each driving
function determines its initial hull and trace.  Localizing before collision
and before the chart derivatives become unbounded, the stochastic integral
in (L8a) is a local martingale.  Its remaining continuous finite-variation
term must therefore vanish identically.  As signed measures,
\[
 h_t''(W_t)\bigl(dC_t-6\,dt\bigr)=0.
 \tag{L8b}
\]
The second derivative starts nonzero by (L8) and remains so on a positive,
possibly random, initial time interval.  Equation (L8b) proves
$C_t=6t$ on that interval.

This conclusion applies at the start of every deterministic marked domain.
To propagate it without any density assumption, set
\[
 T=\inf\{t\geq0:C_t-6t\ne0\},\qquad \inf\varnothing=\infty.
\]
In the usual right-continuous filtration this is a stopping time: the first
entrance time of the continuous adapted process $C_t-6t$ into the open set
$\mathbb R\setminus\{0\}$.  The positive initial zero interval shows
$T>0$ almost surely.  If $T<\infty$, continuity implies $C_T-6T=0$, and by
the definition of $T$ there must be nonzero values arbitrarily soon after
$T$.  Restarting the exploration at $T$ using the strong domain Markov
property gives a positive interval on which
$C_{T+s}-C_T=6s$, a contradiction.  (One may localize $T$ before deterministic
bounds and before any terminal time.)  Hence $C_t=6t$ throughout.
L\'evy's characterization now gives the asserted Brownian driving law,
and the Loewner correspondence identifies the trace as chordal
$\mathrm{SLE}_6$.
\end{proof}

Translation invariance was not needed in this conditional theorem.
The additional condition is substantial: it asks for local martingales
\emph{in the conformal charts of all physical domains}, not just one chosen
half-plane.  A possible route toward universality would be to derive such a
martingale condition from suitable physical observables.  Nothing in this
calculation constructs those observables.

\subsection{A constant-characteristic version and an explicit two-slit test}

A related sufficient assumption is that every marked-domain driving process
has the same deterministic constant variance $\kappa$ and drift $b$, in
its capacity chart.  Under locality (L6), a deletion with initial derivatives
$d=h_0'(0)$ and $e=h_0''(0)$ must obey
\[
 b=\frac b d+\left(\frac\kappa2-3\right)\frac e{d^2}.
 \tag{L9}
\]
Choose the two vertical slits in (L7) at $x$ and $-x$ with the same height.
Their derivatives have the same $d\in(0,1)$ and opposite nonzero $e$.
Subtracting their instances of (L9) yields $\kappa=6$; either equation then
gives $b=b/d$, hence $b=0$.  Thus the only constant-characteristic ensemble
in this class is $\mathrm{SLE}_6$.

This two-slit computation is a useful check on the signs and constants,
but its universality premise is much stronger than physical translation
invariance.  The current bottleneck is to prove an adequate chart-martingale
or characteristic condition from genuinely coarse properties of critical
percolation, rather than to manipulate (L6) further.

\section{Dilation symmetry and criticality}
\subsection{Dilation covariance and subsequential limits}

\begin{lemma}
\label{lem:am-scaleaffine}
Every family with a fixed affine SLE$_6$ normalization is covariant under
all scalar dilations:
$\mathcal P_{rD,rp,rq}=r_*\mathcal P_{D,p,q}$ for $r>0$.
\end{lemma}
\begin{proof}
Dilations commute with the linear part of the normalizer; conjugating by
its full affine map gives a Euclidean similarity. Use SLE similarity covariance.
\end{proof}

\begin{lemma}[The cluster set, not each cluster point, inherits dilations]
\label{lem:am-clusters}
Let $X$ be a topological space of ensembles with continuous dilation
actions $T_r$. Suppose $P_\delta\in X$ has exact mesh covariance
$T_rP_\delta=P_{r\delta}$ and is relatively compact as $\delta\downarrow0$.
Its cluster set $\mathcal C$ satisfies $T_r\mathcal C=\mathcal C$.
If $\mathcal C=\{P\}$, then $T_rP=P$ for every $r>0$.
\end{lemma}
\begin{proof}
If $P_{\delta_n}\to P$, continuity gives
$P_{r\delta_n}=T_rP_{\delta_n}\to T_rP$.
Thus $T_r\mathcal C\subset\mathcal C$; use $T_{1/r}$ for equality.
\end{proof}
For laws indexed by marked domains the action is
$(T_rP)_{D,p,q}=r_*P_{r^{-1}D,r^{-1}p,r^{-1}q}$.
One must choose a topology in which the actions are continuous; this is
not silently supplied by abstract tightness. The distinction is real:
on $X=S^1$, take $P_\delta=\exp(i\omega\log\delta)$ and
$T_rz=\exp(i\omega\log r)z$, $\omega\ne0$.
The cluster set is all of $S^1$ but has no member fixed by every dilation.
Uniform nondegenerate observable bounds do not remove this possibility:
$\tfrac12+\varepsilon\sin(\omega\log(R/\delta))$, $0<\varepsilon<1/2$,
remains uniformly bounded while oscillating in scale phase.
This is a topological example, not a percolation counterexample.
It nevertheless explains why deriving dilation invariance of every
subsequential limit from tightness or RSW alone would be circular.

\subsection{All-scale RSW and a nonzero near-critical parameter}

\begin{proposition}[A finite correlation scale conflicts with all-scale RSW]
\label{prop:am-rswmass}
Suppose continuum colorings $P^\lambda$ obey
$T_rP^\lambda=P^{r^{-\alpha}\lambda}$ for some $\alpha>0$.
For a fixed rectangle shape let $c(\lambda)$ be its unit-scale white
crossing probability. Assume $c(\lambda)\to1$ as $\lambda\to+\infty$ and
$c(\lambda)\to0$ as $\lambda\to-\infty$.
If one $P^\lambda$ has white crossing probabilities in
$[\varepsilon,1-\varepsilon]$ at every scale $R>0$, with fixed
$\varepsilon>0$, then $\lambda=0$.
\end{proposition}
\begin{proof}
Scaling a size-$R$ rectangle down to unit size changes the parameter to
$R^\alpha\lambda$. Its crossing probability is $c(R^\alpha\lambda)$,
which tends to $1$ or $0$ for positive or negative $\lambda$.
\end{proof}

This is a conditional calculation, separate from existence and regularity
of near-critical continuum percolation. For standard triangular-lattice
near-critical percolation $\alpha=3/4$.
RSW merely on each bounded scale range, with constants depending on that
range, does not exclude nonzero mass. Moreover the target-tangent
construction has bounded anisotropy and no preferred scale; removing a
near-critical parameter does not address the target-consistency obstruction.

% Ancillary research fragment. Main target-germ construction is in target_germ.tex.
\section{Near-critical obstruction: what is rigorous and what is not yet transferred}
\label{sec:nearcritical-caveat}

The missing criticality assumption is a separate obstruction from terminal
germ dependence.  On triangular mesh $\eta$, write
\[
 r(\eta)=\frac{\eta^2}{\alpha_4^{\eta}(\eta,1)},\qquad
 p_\eta=\frac12+\lambda r(\eta),\qquad \lambda\ne0,
\]
where $\alpha_4$ is the critical alternating four-arm probability.  The
normalization is $r(\eta)=\eta^{3/4+o(1)}$.  At every mesh the exploration has
exact spatial domain Markov and locality properties for every $p_\eta$.
These two properties alone cannot express criticality.

\paragraph{Established external facts (not reproved here).}
Nolin--Werner, \emph{Asymmetry of near-critical percolation interfaces},
JAMS 22 (2009), 797--819,
\url{https://arxiv.org/abs/0710.1470}, establish subsequential limits of
near-critical interfaces and their singularity with respect to ordinary
$\mathrm{SLE}_6$.  Garban--Pete--Schramm, \emph{The scaling limits of
near-critical and dynamical percolation}, JEMS 20 (2018), 1195--1268,
\url{https://arxiv.org/abs/1305.5526}, construct the configuration-valued
near-critical scaling limit and establish its conformal covariance.  Their
Section 13.2 explicitly distinguishes the measurable trace of an interface
from its order as a curve: the latter uniqueness is not proved there.
Their Section 11 Markov property concerns change of the near-critical
parameter, not the spatial exploration stopping-time property needed here.
Thus neither statement may silently be substituted for full convergence of
ordered exploration kernels in random slit domains.  A 2025 Durham thesis
of L\'eonie Papon, introduction,
\url{https://etheses.durham.ac.uk/id/eprint/16026/1/Papon-thesis-v2.pdf},
also records ordered near-critical percolation interface uniqueness as a
remaining challenge.  These observations are limits of the sources checked,
not a proof that no later theorem exists.

\paragraph{A simple necessary test for fixed affine SLE.}
Every fixed affine image of the ordinary SLE ensemble is covariant under
all scalar dilations.  Indeed, if the linear part of the affine map is $L$,
then $L(rx)=rLx$; the translation part is absorbed by translation covariance.
Consequently a genuine failure of dilation covariance excludes \emph{every}
fixed affine normalization.  Singularity relative to ordinary SLE alone does
not have this implication: an anisotropic affine image may also be singular.

Under a physical dilation by $s>0$, the near-critical parameter changes as
\[
 s\,\omega^{\lambda}\ \stackrel{\mathrm{law}}{=}
 \omega^{s^{-3/4}\lambda}.
\]
The sign of the exponent follows already on the lattice: the original mesh
$\eta$ becomes $s\eta$ while the same occupation probability has effective
parameter $\lambda r(\eta)/r(s\eta)\to s^{-3/4}\lambda$.  Known nonconstant
crossing probabilities as functions of $\lambda$ therefore witness a
macroscopic scale in the configuration law.  Transferring this witness to
a full domain-indexed curve ensemble requires the ordered-kernel step just
identified.  This investigation therefore records near-critical interfaces
as a \emph{conditional additional counterexample}, not as the proof of the
counterexample in the main text.

\subsection{An explicit limit-transfer lemma}

Here is a precise sufficient condition for the commonly asserted passage of
the spatial Markov property to scaling limits.  It isolates the missing
assumption rather than labeling it automatic.

Let $(I_n,D_n,Y_n)$ denote respectively the explored initial segment, remaining
marked domain, and future curve at a stopping time.  Suppose these random
objects take values in Polish spaces and converge jointly in distribution
to $(I,D,Y)$.  Let $K_n$ be the future-curve kernels, so
\[
 \mathbb E[f(Y_n)\mid I_n,D_n]=K_n f(D_n).
\]
Assume there is a Borel probability kernel $K$ such that, for every bounded
continuous $f$ and every sequence $d_n\to d$ with $d$ in a full-law set for
$D$, one has $K_n f(d_n)\to Kf(d)$.  Then
\[
 \mathbb E[f(Y)\mid I,D]=Kf(D).
\]

\begin{proof}
By a Skorokhod representation the triples may be assumed to converge almost
surely.  For any bounded continuous test function $h$ of $(I,D)$, the exact
finite-$n$ identity gives
\[
 \mathbb E[h(I_n,D_n)f(Y_n)]
 =\mathbb E[h(I_n,D_n)K_n f(D_n)].
\]
Bounded convergence and the sequential-kernel hypothesis give the same
identity at the limit, with $Kf(D)$ on the right.  The monotone class theorem
extends the identity to bounded measurable $h$, and a countable
convergence-determining class of $f$ identifies the conditional kernel.
\end{proof}

The decisive hypothesis controls kernels at \emph{random slit domains};
convergence for one deterministic Jordan domain is insufficient.  Analogous
coupled convergence can transfer locality before first approach to a removed
closed set: stop first at positive-distance neighborhoods, where uniform
curve convergence leaves the two coupled initial segments identical, then
decrease the distance to zero.  Exact first-hit stopping needs a separate
continuity argument and is not implicit in this observation.


% Research fragment; intended for inclusion in progress.tex.
\section{Affine identifiability without Jordan boundary regularity}
\label{sec:rough-affine}

\paragraph{Status and interpretation.}
The extension proved here uses only standard quasiconformal compactness,
prime-end extension for quasiconformal disk maps, Fatou's theorem for bounded
analytic functions, and uniqueness of bounded harmonic functions from their
almost-everywhere disk boundary values.  It removes the Jordan assumption
from affine metric identifiability.  At completely arbitrary prime ends a
conformal image of SLE need not have a Euclidean-continuous endpoint; its
law is therefore understood in the prime-end compactification, or as its
conformal hull exploration.  The result applies in particular whenever the
physical curves are continuous, as in the exploration setting of interest.
No classification of arbitrary exploration ensembles is asserted.

\subsection{A quasiconformal disk map fixing the boundary stays close to the identity}

\paragraph{Lemma (proved).}
For each finite $K\geq1$ there exists $C(K)<\infty$ such that every
orientation-preserving $K$-quasiconformal self-homeomorphism
$h:\mathbb D\to\mathbb D$ whose boundary extension fixes every point of
$\partial\mathbb D$ satisfies
\[
 d_{\mathbb D}(h(z),z)\leq C(K)\qquad(z\in\mathbb D).
 \tag{RA1}
\]
Here $d_{\mathbb D}$ is hyperbolic distance.

\begin{proof}
First there is $r(K)<1$ with $|h(0)|\leq r(K)$ for all such maps.  To see
this, extend $h$ to the sphere by reflection in the unit circle:
\[
 \widetilde h(z)=\frac1{\overline{h(1/\bar z)}}\quad(|z|>1).
\]
This is a $K$-quasiconformal sphere map, and fixes three prescribed distinct
points of the unit circle.  The standard compactness theorem for normalized
$K$-quasiconformal sphere maps, including compactness of their inverses,
precludes a sequence with $h_n(0)\to\zeta\in\partial\mathbb D$: a limiting
quasiconformal homeomorphism would send both $0$ and the fixed boundary
point $\zeta$ to $\zeta$.  Thus such an $r(K)$ exists.
For an arbitrary $z$ choose a disk automorphism $a_z$ sending $z$ to $0$.
The map $a_z\circ h\circ a_z^{-1}$ is again $K$-quasiconformal with identity
boundary map, so
\[
 |a_z(h(z))|\leq r(K).
\]
Conformal invariance of the hyperbolic metric gives (RA1), for example with
$C(K)=2\operatorname{arctanh}r(K)$ in the usual normalization.
\end{proof}

\paragraph{Consequence (proved).}
For every $\zeta\in\partial\mathbb D$, the image of the radial approach
$r\zeta$ under $h$ is a nontangential approach to the same $\zeta$, with
aperture bounded solely in terms of $K$.
Indeed rotate so $\zeta=1$ and write
\[
 h(r)=\frac{r+u_r}{1+ru_r},\qquad |u_r|\leq r(K).
\]
Then
\[
 |1-h(r)|\leq(1-r)\frac{1+r(K)}{1-r(K)},
\]
while
\[
 1-|h(r)|\geq
 \frac{1-r}{2}\frac{1-r(K)}{1+r(K)}.
\]
The ratio $|1-h(r)|/(1-|h(r)|)$ is therefore uniformly bounded.  In
particular $h(r\zeta)\to\zeta$ within a fixed Stolz region.

\subsection{The identifiability theorem on any bounded simply connected domain}

\paragraph{Theorem (proved).}
Let $D$ be any bounded simply connected planar domain and let $p,q$ be
distinct prime ends.  For orientation-preserving invertible real-linear
maps $A,B$, define
\[
 \nu^A_{D,p,q}
 =(A^{-1})_*\operatorname{Law}(\mathrm{SLE}_6(AD;Ap,Aq)).
\]
Then
\[
 \nu^A_{D,p,q}=\nu^B_{D,p,q}
 \quad\Longleftrightarrow\quad
 BA^{-1}\text{ is a Euclidean similarity}
 \quad\Longleftrightarrow\quad
 B^{\mathsf T}B=cA^{\mathsf T}A\quad(c>0).
 \tag{RA2}
\]
The statement holds for exploration laws in the prime-end compactification,
and hence for their physical curves whenever defined continuously.

\begin{proof}
Set $M=BA^{-1}:AD\to BD$.  This bilipschitz map carries null crosscut
chains to null crosscut chains, with equivalent chains remaining equivalent;
thus it induces the natural homeomorphism of prime-end compactifications.
Choose conformal charts $\phi_A:AD\to\mathbb H$ and
$\phi_B:BD\to\mathbb H$ sending the corresponding source and target to
$0,\infty$.  The interior map
\[
 H=\phi_B\circ M\circ\phi_A^{-1}:\mathbb H\to\mathbb H
\]
is quasiconformal with a boundary extension to $\mathbb R\cup\{\infty\}$,
fixing $0,\infty$.  The boundary map is increasing.  The boundary-swallowing
race of the main record remains a prime-end topological event: an affine
homeomorphism preserves the order in which boundary ends are cut off from
the target by the explored filled hull.  Equality of the two path laws,
followed by strictness of that race probability, therefore gives, for every
$x<0<y$,
\[
 \frac{-H(x)}{H(y)-H(x)}=\frac{-x}{y-x}.
\]
Consequently $H(x)=cx$ on the entire extended real boundary, for some $c>0$.
This step uses prime-end boundary values of $H$, not Euclidean boundary
values of either Riemann map.

Pass to disk charts, absorbing the conformal map $z\mapsto cz$ into the
second chart.  We obtain bounded conformal maps
\[
 \psi_A:\mathbb D\to AD,\qquad \psi_B:\mathbb D\to BD
\]
such that
\[
 h=\psi_B^{-1}\circ M\circ\psi_A
\]
is a $K$-quasiconformal disk self-homeomorphism with identity boundary map.
Thus pointwise in the disk,
\[
 M\psi_A=\psi_B\circ h.
 \tag{RA3}
\]
Fatou's theorem gives radial limits $\psi_A^*(\zeta)$ and nontangential
limits $\psi_B^*(\zeta)$ for almost every circle point $\zeta$.
The preceding lemma shows that $h(r\zeta)$ approaches that same $\zeta$
nontangentially.  Taking $r\uparrow1$ in (RA3) yields
\[
 M\psi_A^*(\zeta)=\psi_B^*(\zeta)
 \quad\text{for almost every }\zeta\in\partial\mathbb D.
 \tag{RA4}
\]
Both real coordinates of $M\psi_A$ are bounded harmonic functions on the
disk, since $M$ is real-linear and $\psi_A$ is bounded analytic.
The coordinates of $\psi_B$ are also bounded harmonic.  Uniqueness of the
Poisson representation for bounded harmonic functions now gives
\[
 M\psi_A=\psi_B\qquad\text{throughout }\mathbb D.
\]
Thus $M$ is conformal on the nonempty open set $AD$.  Writing
$M(z)=\alpha z+\beta\bar z$, holomorphicity of $M\psi_A$ gives
$\beta\overline{\psi_A'(z)}=0$.  Since $\psi_A'$ never vanishes,
$\beta=0$.  Invertibility and orientation imply that $M$ is a direct
Euclidean similarity.  The converse is SLE similarity covariance.
Reflections reduce orientation-reversing normalizers to this case.
\end{proof}

The boundedness assumption was used exactly when invoking bounded analytic
Fatou limits and bounded harmonic uniqueness.  No assertion about arbitrary
unbounded domains follows without a growth or uniqueness condition.

\subsection{A necessary and sufficient DMP condition inside the affine ansatz}

Suppose, as an explicit additional ansatz, that an ensemble has the form
\[
 P_{D,p,q}=\nu^{A(D,p,q)}_{D,p,q}
 \tag{RA5}
\]
for a deterministic measurable affine label on each marked bounded simply
connected domain.  Remove the inessential positive scalar from its metric by
putting
\[
 T(D,p,q)=
 \frac{A(D,p,q)^{\mathsf T}A(D,p,q)}
 {\sqrt{\det(A(D,p,q)^{\mathsf T}A(D,p,q))}}.
 \tag{RA6}
\]
Thus $T$ is a positive-definite symmetric matrix of determinant one.
Assume the usual Borel family and strong conditional-kernel formulation, so
that equality of conditional path laws at a stopping time is an equality of
probability kernels almost surely.

\paragraph{Corollary (proved).}
Within ansatz (RA5), the strong domain Markov property holds if and only if,
for every starting marked domain and every admissible stopping time $\tau$,
\[
 T(D_\tau,p_\tau,q)=T(D,p,q)\qquad\text{almost surely}.
 \tag{RA7}
\]
Here $D_\tau$ is the target component of the remaining domain.  It need not
have Jordan boundary.

\begin{proof}
Fix the starting label $A=A(D,p,q)$.  Under the fixed-$A$ affine SLE law,
ordinary SLE's strong domain Markov property gives, conditional on the
explored initial segment, the remaining law
\[
 \nu^A_{D_\tau,p_\tau,q}.
\]
Strong DMP for (RA5) demands that this be
$\nu^{A(D_\tau,p_\tau,q)}_{D_\tau,p_\tau,q}$ almost surely.
The rough-domain theorem (RA2), applied at each realized slit domain,
is equivalent to equality of their normalized metric tensors, namely
(RA7).  Conversely (RA7) and similarity covariance identify these two
conditional laws, proving strong DMP.
\end{proof}

This corollary classifies the DMP constraint only \emph{within the stated
ansatz}: the affine tensor must be a conserved quantity of the exploration
state.  A target-germ tensor supplies exactly such a conserved quantity.
The corollary neither proves that a general ensemble has affine-SLE form nor
that every conserved tensor must arise from a target germ.  Locality,
translation covariance, target consistency, and boundary continuity impose
further constraints which are not characterized by (RA7) alone.

% Independent verification fragment: transfer of local target consistency.
\section{Local target consistency passes to marginal curve limits}
\label{sec:local-target-transfer}

\paragraph{Scope.}
This lemma transfers an exploration coupling only inside a fixed neighborhood
of the common start. It needs convergence of the two marginal curve laws, but
neither convergence of random slit-domain kernels nor a limiting strong domain
Markov theorem. It does not assert a coupling of the limiting curves all the
way to a distant boundary arc.

\subsection{Stopping an ordered unparameterized curve on a circle}

Fix $p\in\mathbb R^2$ and $R_0>0$. Work in the usual metric space of oriented
continuous curves modulo increasing reparameterization, with the uniform
Fr\'echet topology. Curves start at $p$ and end outside
$\overline{B(p,R_0)}$. Small changes of the starting points are harmless;
for the sequence version they need only lie inside every fixed smaller ball
for sufficiently large indices.
For a continuous representative $\gamma:[0,1]\to\mathbb R^2$ and
$0<r<R_0$, put
\[
 \tau_\gamma(r)=\inf\{t:|\gamma(t)-p|\geq r\},\qquad
 E_r\gamma=\gamma|_{[0,\tau_\gamma(r)]}.
 \tag{LT1}
\]
The stopped path is regarded as an ordered unparameterized curve, so this
definition is independent of pauses or changes of clock.

\begin{lemma}[Almost every circle is a continuity circle]
\label{lem:lt-stopcontinuity}
For each fixed continuous curve $\gamma$, the map
$(\eta,s)\mapsto E_s\eta$ is continuous at $(\gamma,r)$ for all but at most
countably many radii $r\in(0,R_0)$. The stopping map is Borel.
Consequently, for any probability law $\mu$ on these curves, $E_r$ is
$\mu$-almost surely continuous for Lebesgue-almost every deterministic $r$.
\end{lemma}
\begin{proof}
Choose any continuous representative of $\gamma$.
The function $r\mapsto\tau_\gamma(r)$ is nondecreasing and hence has at
most countably many discontinuities. Suppose it is continuous at $r$.
If $(\gamma_n,r_n)\to(\gamma,r)$ in the product of Fr\'echet and ordinary
metrics, choose increasing reparameterizations for which representatives
$\widetilde\gamma_n$ converge uniformly to $\gamma$.
Writing $\varepsilon_n=\|\widetilde\gamma_n-\gamma\|_\infty+|r_n-r|$,
first-hit times obey, for every slightly larger positive error tending to zero,
\[
 \tau_\gamma(r-2\varepsilon_n)
 \ \leq\ \tau_{\widetilde\gamma_n}(r_n)
 \ \leq\ \tau_\gamma(r+2\varepsilon_n).
 \tag{LT2}
\]
If $\varepsilon_n=0$ the assertion is immediate; using a strictly larger error
avoids any immaterial equality convention. Continuity of $\tau_\gamma$ at
$r$ gives convergence of the middle times to $\tau_\gamma(r)$.
Uniform continuity of $\gamma$ then makes the stopped representatives,
linearly recoded onto $[0,1]$, converge uniformly. This proves joint
continuity at $(\gamma,r)$.

The stopping map is understood on all curves reaching the selected radius, including already stopped curves reaching a larger radius. For parameterized curves the stopping time is Borel, since
\[
 \{\tau_\gamma(r)\leq t\}
 =\left\{\max_{0\leq u\leq t}|\gamma(u)-p|\geq r\right\}.
\]
Stopping and linear recoding are therefore Borel path operations. They
commute with reparameterization and give the usual Borel stopping map on
ordered unparameterized curves.

For completeness, the almost-every-radius assertion does not require a
measurable choice of representatives. For a map between metric spaces its
set of continuity points is a $G_\delta$ set, obtained from its oscillations
on neighborhoods. Apply this to the joint stopping map. Its discontinuity
set in the product of curve space and radius space is Borel, and every
vertical section has Lebesgue measure zero by the countability argument
above. Fubini's theorem now gives the assertion under $\mu$.
\end{proof}

\paragraph{Why a bad radius can occur.}
A path may first touch a circle, then spend a nontrivial segment inside it,
and only later cross it. A small outward perturbation can postpone its
first exit until that later crossing. This is why stopping continuity
cannot simply be claimed at every deterministic radius for every path.
The preceding lemma handles these tangencies without excluding them by an
unproved regularity assumption.

\subsection{Transfer using only two marginal limits}

\begin{theorem}[Transfer of local target consistency]
\label{thm:lt-transfer}
For $i=1,2$, let $\gamma_n^i$ be oriented continuous random curves whose
marginal laws converge weakly in the uniform Fr\'echet topology to laws
$\mu^i$. Assume the limiting curves start at $p$ and end outside
$\overline{B(p,R_0)}$. Suppose that for every fixed $r\in(0,R_0)$ and all
sufficiently large $n$,
\[
 \operatorname{Law}(E_r\gamma_n^1)
 =\operatorname{Law}(E_r\gamma_n^2).
 \tag{LT3}
\]
Then the same equality holds for the limiting curves at \emph{every}
$r\in(0,R_0)$:
\[
 (E_r)_*\mu^1=(E_r)_*\mu^2.
 \tag{LT4}
\]
The hypothesis (LT3) may equivalently be replaced by a discrepancy tending
to zero in any metric inducing weak convergence of the stopped laws.
\end{theorem}
\begin{proof}
Lemma~\ref{lem:lt-stopcontinuity} gives a full-Lebesgue-measure set of
radii at which $E_r$ is almost surely continuous under both limiting laws.
At each such radius the continuous mapping theorem and (LT3) give (LT4).
This part needs no joint convergence of $(\gamma_n^1,\gamma_n^2)$.

To obtain every smaller radius, fix arbitrary $r<R_0$ and choose a good
$s\in(r,R_0)$. Equality of the stopped-at-$s$ laws remains equality after
applying the Borel map $E_r$. Since
$E_r(E_s\gamma)=E_r\gamma$ for $r<s$, this proves (LT4) at $r$ without
requiring any continuity of $E_r$ there.
\end{proof}

\subsection{Application to two percolation targets}

Let $D$ be a Jordan domain, $p$ a common boundary start, and let the targets
$q_1,q_2$ lie on a boundary arc $J$ with $p\notin\overline J$.
The boundary colors for the two Dobrushin explorations differ only on the
arc between the targets. Use identical independent bulk colors and identical
boundary colors outside a lattice approximation $J_n$ of this disagreement
arc, with a common discrete starting edge. The discrete exploration rules
produce exactly the same ordered initial path until the exploration first
encounters $J_n$.
Choose
\[
 0<R_0<\operatorname{dist}(p,\overline J).
\]
For all large $n$, the approximating disagreement arc stays outside the
closed ball $\overline{B(p,R_0)}$. Thus (LT3) holds inside every smaller
ball. If a common subsequence makes each marginal exploration converge,
Theorem~\ref{thm:lt-transfer} supplies local target consistency for their
limiting initial segments. The lattice approximations must in fact have the
stated common-start/common-color coupling; a small displacement of two
unrelated explorations alone would not justify pathwise closeness.

\paragraph{Consequence for affine decorations.}
The affine-metric elimination argument needs only equality of these local
prefix laws, not equality until the full distant arc. Suppose $D$ is
analytic near $p$ and the two limiting marginals have laws
\[
 (A_i^{-1})_*\operatorname{SLE}_6(A_iD,A_ip,A_iq_i),\qquad i=1,2.
\]
Translate $p$ to zero and blow up space. By conformal kernel convergence at
the smooth start, the local initial laws tend to
\[
 (A_i^{-1})_*\operatorname{SLE}_6(A_iH_p,0,\infty),
\]
where $H_p$ is the tangent half-plane. Equality in every sufficiently small
physical ball transfers after dilation to equality of the limiting stopped
laws in every fixed half-disk (first at common continuity radii, then all
smaller radii by the same composition argument). These local restrictions
determine the entire half-plane law. Strict left-passage identifiability
therefore makes $A_2A_1^{-1}$ a similarity.

Consequently the target-tangent affine counterexample cannot be the family
of simultaneous two-target subsequential percolation limits on a fixed
smooth domain with targets carrying different tangent directions. This
exclusion is independent of the more difficult transfer of a strong domain
Markov property to random limiting slit domains.

% Independent research fragment: a rigorous partial result, not a counterexample.
% Independent research fragment: a rigorous partial result, not a counterexample.
\section{Affine metrics encoded in initial path germs}
\label{sec:source-germ-decoder}

\paragraph{Status and motivation.}
The result below is proved for the \emph{oriented initial trace germ} of an
affine image of half-plane SLE$_6$ at its smooth starting boundary point.
It does not construct a source-dependent domain Markov family. In particular,
it does not establish that the metric can be recovered from the unlabelled
boundary germ at a later random exploration tip. Those distinctions are
important when assessing whether target consistency alone could repair the
classification conjecture.

\subsection{State space and the meaning of a trace germ}

We retain the chronological order of a path but forget its original clock.
For the half-plane curves used here, the usual half-plane-capacity clock
provides a canonical representative: the filled hull at time $t$ has
hydrodynamic map $g_t(z)=z+2t/z+o(1/z)$.
This clock is determined by the ordered unparameterized trace. Section~\ref{sec:loewner-admissible} justifies continuous Loewner driving of affine images; hence they admit this clock.
The clock has infinite total lifetime almost surely: its total half-plane
capacity $T$ is positive, while spatial scale invariance makes its law equal
to that of $r^2T$ for every $r>0$; the only such probability laws on
$[0,\infty]$ are supported on $\{0,\infty\}$. Thus $T=\infty$.
All functions below are functions of the ordered unparameterized trace.

Use the standard Borel space of these capacity-parameterized continuous
curves, with the sigma-field generated by restrictions to finite time
intervals. Equivalently, this can be regarded as a Borel realization of the
usual ordered unparameterized path space. For $t>0$, let
\[
 \mathcal F_t=\sigma\{\gamma(s):0\leq s\leq t\},\qquad
 \mathcal F_{0+}=\bigcap_{t>0}\mathcal F_t.
\]
The latter is the initial trace-germ sigma-field. It is unaffected by changing
any part of the curve after a positive capacity time. The canonical time is
used only to give an explicit measurable coding of this germ.

For $r>0$, let $S_r\gamma$ be the spatially dilated trace, again put in the
canonical clock. Capacity scaling gives the exact identity
\[
 (S_r\gamma)(t)=r\gamma(t/r^2).
 \tag{SG1}
\]
In particular $S_r$ is invertible with inverse $S_{1/r}$.

\subsection{Dilation ergodicity}

\begin{lemma}[Brownian scaling is mixing]
\label{lem:sg-brownianmix}
For a one-sided Brownian motion $B=(B_t)_{t\geq0}$ and any fixed $r>1$, the
measure-preserving transformation
\[
 (\mathcal S_rB)_t=rB_{t/r^2}
\]
is mixing and hence ergodic.
\end{lemma}
\begin{proof}
Set $Z_s=e^{-s/2}B_{e^s}$ for $s\in\mathbb R$.
This is a stationary centered Gaussian process with
\[
 \mathbb E[Z_sZ_u]=e^{-|s-u|/2}.
\]
It is the stationary Ornstein--Uhlenbeck process. Under $\mathcal S_r$, it
becomes $Z_{s-2\log r}$. Covariances between finite collections of coordinates
separated by $n(2\log r)$ tend to zero, so the corresponding jointly Gaussian
vectors become independent in the limit. Approximation by finite-coordinate
cylinder functions gives mixing on the full path sigma-field.
The correspondence is invertible almost surely:
$B_t=\sqrt t\,Z_{\log t}$ for $t>0$, with $B_0=0$. Thus mixing transfers back
to Brownian scaling. The two-sided log-time process here is obtained from a
one-sided Brownian path at times $t>0$; no independent Brownian past is being
added.
\end{proof}

\begin{lemma}[Affine SLE$_6$ is dilation mixing]
\label{lem:sg-affinemix}
Let $A$ be any invertible linear map carrying $\mathbb H$ onto itself and
fixing its positive boundary orientation. Define
\[
 \mu_A=\operatorname{Law}\bigl(A^{-1}\eta\bigr),\qquad
 \eta\sim\operatorname{SLE}_6(\mathbb H,0,\infty),
\]
with the physical half-plane-capacity representative understood.
For every fixed $r>1$, $S_r$ preserves $\mu_A$ and is mixing under $\mu_A$.
\end{lemma}
\begin{proof}
The SLE trace is a measurable, scaling-equivariant function of its Brownian
Loewner driver $\sqrt6B$. The trace exists almost surely. Therefore its
dilation action is a factor of the Brownian scaling action and is mixing
by Lemma~\ref{lem:sg-brownianmix}. The map on ordered traces
$\eta\mapsto A^{-1}\eta$ commutes with scalar spatial dilations, and is an
invertible measurable map after recoding by capacity. Mixing and invariance
therefore transfer to $\mu_A$.
\end{proof}

\subsection{Separating distinct metrics from arbitrarily short initial segments}

Normalize the maps preserving $\mathbb H$ by
\[
 A_{b,c}=\begin{pmatrix}1&b\\0&c\end{pmatrix},
 \qquad b\in\mathbb R,\ c>0.
 \tag{SG2}
\]
This parametrizes all constant positive-definite conformal metrics, up to a
positive scalar, in a fixed half-plane frame: the tensor is
$A_{b,c}^{\mathsf T}A_{b,c}$ with its $(1,1)$ entry normalized to $1$.
The laws $\mu_{A_{b,c}}$ are distinct. Indeed, the strictly increasing SLE$_6$
left-passage function $F$ gives
\[
 \mu_{A_{b,c}}[L(x+iy)]=F\left(\frac{x+by}{cy}\right).
 \tag{SG3}
\]
Equality of the probabilities at $z=i$ and $z=1+i$ first identifies $b/c$
and then $1/c$, hence both parameters.

\begin{lemma}[Unit-time restrictions distinguish the metrics]
\label{lem:sg-unitrestriction}
If $A\ne A'$ are two normalized maps in (SG2), the laws of their traces
restricted to capacity time $[0,1]$ are distinct.
\end{lemma}
\begin{proof}
If the restrictions agreed, dilation invariance and (SG1) would make the
restrictions to $[0,T]$ agree for every $T>0$. Those restrictions generate the
full path sigma-field, so the entire laws would agree, contrary to (SG3).
\end{proof}

\begin{theorem}[Metric decoding from an initial trace germ]
\label{thm:sg-decoder}
Let $\{A_j:j\in J\}$ be any finite or countable collection of distinct
normalized maps (SG2). There exists an $\mathcal F_{0+}$-measurable function
$D$ with values in $J\cup\{\dagger\}$ such that
\[
 \mu_{A_j}[D(\gamma)=j]=1\qquad\hbox{for every }j\in J.
 \tag{SG4}
\]
In particular, restrictions of these affine SLE$_6$ measures to the initial
trace-germ sigma-field are mutually singular, even though each separately
has its usual almost-sure germ zero-one property.
\end{theorem}
\begin{proof}
For each pair $j\ne k$, Lemma~\ref{lem:sg-unitrestriction} supplies a Borel
set of unit-time traces whose indicator $f_{jk}$ has distinct expectations
under $\mu_{A_j}$ and $\mu_{A_k}$. There are only countably many such pairs;
list all the resulting bounded unit-time observables as $(f_m)_{m\geq1}$.
Write $v_{j,m}=\int f_m\,d\mu_{A_j}$. The vectors
$v_j=(v_{j,m})_m$ are pairwise distinct.

Fix $r>1$. By Lemma~\ref{lem:sg-affinemix} and Birkhoff's ergodic theorem,
for each $m$ and each $j$,
\[
 \frac1N\sum_{n=1}^N f_m(S_{r^n}\gamma)
 \longrightarrow v_{j,m}\qquad\mu_{A_j}\hbox{-almost surely}.
 \tag{SG5}
\]
The intersection of these events over $m$ still has full $\mu_{A_j}$ measure.
Every term $f_m(S_{r^n}\gamma)$ depends only on the trace up to time
$r^{-2n}$. For any $t>0$, all sufficiently late terms are consequently
$\mathcal F_t$-measurable. Omitting finitely many initial terms does not
change a Ces\`aro limit. Thus every limit in (SG5), and the event that the
limit exists, is $\mathcal F_{0+}$-measurable.

Define $D=j$ if all limits exist and their vector equals $v_j$; use
$D=\dagger$ otherwise. Countability of $J$ and pairwise distinctness of the
$v_j$ make this a measurable unambiguous definition. Formula (SG5) proves
(SG4).
\end{proof}


\subsection{A common decoder for the continuum of normalized metrics}

The countable proof extends to all parameters $(b,c)\in\Theta
=\mathbb R\times(0,\infty)$, provided the standard path coding is used
jointly in the affine parameter. We state this coding point explicitly.

\paragraph{Joint Borel capacity coding.}
Let $\eta$ be an ordinary capacity-parameterized SLE$_6$ trace and $K_s$ its
filled hull. Before recoding, its affine image is the continuous path
$A_{b,c}^{-1}\eta(s)$. Its physical capacity clock is
\[
 u_{b,c,\eta}(s)=\tfrac12\operatorname{hcap}(A_{b,c}^{-1}K_s).
 \tag{SG6}
\]
Use the standard measurable coding of compact half-plane hulls, in which
filling a continuous path prefix and taking its half-plane capacity are
Borel operations. Linear transformation is jointly continuous on compact
sets. Consequently $(b,c,\eta,s)\mapsto u_{b,c,\eta}(s)$ is Borel.
The inverse clock is measurable because, on the set where the clock is
continuous and strictly increasing,
\[
 u_{b,c,\eta}^{-1}(t)=\inf\{s:u_{b,c,\eta}(s)>t\}
\]
can be computed using rational $s$ (with equality interpreted by continuity).
Thus evaluation of the recoded image on $[0,1]$ is jointly Borel in
$(b,c,\eta)$. The good-clock conditions are Borel using rational tests;
on their complement use a cemetery path. For each fixed $(b,c)$ this
complement has SLE probability zero by the affine local-growth and lifetime
arguments above. No single full-probability set valid simultaneously for
all affine parameters is required.

This paragraph uses the usual deterministic measurability of hull capacity
as part of the standard Borel curve coding. If that coding fact is not
taken as established, the next corollary should be read conditionally on
joint Borel recoding. The preceding countable theorem needs only measurable
recoding for each of its countably many fixed maps.

\begin{corollary}[Continuum metric decoder]
\label{cor:sg-continuum}
With the joint Borel coding just specified, there exists one
$\mathcal F_{0+}$-measurable map
\[
 D_{\mathrm{all}}:\mathcal C\longrightarrow\Theta\cup\{\dagger\}
\]
such that for every $(b,c)\in\Theta$,
\[
 \mu_{A_{b,c}}[D_{\mathrm{all}}(\gamma)=(b,c)]=1.
 \tag{SG7}
\]
The same single measurable rule therefore recovers every normalized affine
metric, with almost-sure correctness under each metric law separately.
\end{corollary}
\begin{proof}
Choose a countable generating $\pi$-system for the standard Borel space of
unit-capacity path segments, including the cemetery state. For example, use
finite intersections of rational-time evaluation cylinders in rational
rectangles, together with the cemetery singleton. Let $(f_m)_{m\geq1}$ be
their indicator functions. They determine probability laws on unit segments.
Joint Borel recoding and integration against the fixed ordinary SLE law imply
that
\[
 v:\Theta\to[0,1]^{\mathbb N},\qquad
 v_m(b,c)=\int f_m\,d\mu_{A_{b,c}},
 \tag{SG8}
\]
is Borel. Lemma~\ref{lem:sg-unitrestriction} and the determining property
of $(f_m)$ imply that $v$ is injective.
Both parameter and sequence spaces are standard Borel, so the Lusin--Souslin
theorem makes $v(\Theta)$ Borel and $v^{-1}:v(\Theta)\to\Theta$ Borel.

Fix $r>1$. Form the vector $L(\gamma)$ of Ces\`aro limits
\[
 L_m(\gamma)=\lim_{N\to\infty}\frac1N\sum_{n=1}^N
 f_m(S_{r^n}\gamma),
\]
using a cemetery value if any limit fails to exist. The proof of
Theorem~\ref{thm:sg-decoder} shows that $L$ is
$\mathcal F_{0+}$-measurable. Define
$D_{\mathrm{all}}=v^{-1}(L)$ when $L\in v(\Theta)$ and set it to
$\dagger$ otherwise. This is a measurable definition; no measurable
selection theorem is needed.
For each fixed $(b,c)$, ergodicity and Birkhoff applied to the countable
family $(f_m)$ give $L=v(b,c)$ almost surely under $\mu_{A_{b,c}}$.
This proves (SG7).
\end{proof}

\paragraph{Quantifiers and limitations.}
The statement is one common decoder with correctness under each law:
$\forall(b,c)\ \mu_{A_{b,c}}(D_{\mathrm{all}}=(b,c))=1$.
It does not assert one common full-probability set of Brownian seeds on
which correctness holds simultaneously after every uncountable collection
of affine transformations. More importantly, the decoded object is still
an ordered initial trace germ at a smooth start, not the bare boundary germ
of a random slit-domain tip.

\paragraph{What has and has not been proved.}
The theorem demonstrates an exact mechanism by which a microscopic ordered
path germ can store a constant affine metric. It is more than identifiability
of the law: a single arbitrarily short trace germ distinguishes any prescribed
countable collection of metrics almost surely, and the continuum corollary
extends this under the stated joint Borel coding. The result does not contradict
a germ zero-one law, since that law concerns one probability measure, while
the distinguishing full-measure germ sets differ between measures.

A hypothetical source-germ decoration could choose a metric from the smooth
initial source geometry, then attempt to recover the same metric from the
slit boundary at all later tips. That recovery does \emph{not} follow from Theorem~\ref{thm:sg-decoder}.
Three obstacles remain:
\begin{enumerate}
\item A later tip is not a smooth half-plane starting boundary point; the
nearby past trace can have a different two-sided or conditioned local law.
\item The remaining marked domain gives an embedded boundary germ, which may
lose chronological information when the trace has touched itself. The theorem
uses an oriented initial trace germ and does not prove that the needed data
are recoverable from that boundary alone.
\item A domain Markov construction would need one measurable rule valid at
all relevant stopping times and compatible with locality deletions, not merely
a decoder at a deterministic initial point under each fixed metric law.
\end{enumerate}
Extending local metric decoding to random interior or tip germs, or proving
that such decoding is impossible from the domain state, is a concrete next
research question. No full counterexample satisfying target consistency is
claimed from this mechanism.

\section{A fixed metric follows inside the affine-SLE ansatz with target consistency and reversal}
\label{sec:affine-ansatz-classification}

\paragraph{Status.}
This theorem assumes that every marginal is already an affine image of
SLE$_6$. It does not derive that ansatz from coarse percolation properties.
It shows that quite weak physical coherence conditions then eliminate all
variation of the affine conformal structure between domains and marked points.
Geometric reversal is an explicit additional hypothesis; it is not asserted
here to follow from RSW or from the original three axioms.

Let $\mathcal D$ be the class of bounded Jordan domains whose boundaries
consist of finitely many regular real-analytic arcs meeting at corners.
Each analytic arc extends to a neighborhood of its closed parameter interval.
Restrict marks to regular boundary points, meaning interior points of these
arcs. Suppose
\[
 P_{D,p,q}=(A_{D,p,q}^{-1})_*\operatorname{SLE}_6
 (A_{D,p,q}D,A_{D,p,q}p,A_{D,p,q}q),
 \tag{AA1}
\]
where $A_{D,p,q}$ is an invertible orientation-preserving linear map. Write
$[Q_{D,p,q}]$ for the projective positive-definite tensor
$[A_{D,p,q}^{\mathsf T}A_{D,p,q}]$; brackets identify positive scalar multiples.

The following elementary local identifiability principle will be used:
if two laws of the form (AA1), on domains agreeing near the same regular
start, have identical initial paths stopped in every sufficiently small ball,
then their projective tensors are equal. Indeed blow up at the regular start.
Conformal kernel convergence for ordinary SLE gives half-plane limits
$A_i^{-1}\operatorname{SLE}_6(A_iH,0,\infty)$. Equality of local prefixes
gives equality of these half-plane laws; the strictly monotone SLE$_6$
left-passage function makes $A_2A_1^{-1}$ a similarity. This argument uses
only local SLE convergence and the affine ansatz.

\begin{theorem}[Conditional global metric rigidity]
\label{thm:aa-global}
Assume (AA1), together with the following properties on $\mathcal D$:
\begin{enumerate}
\item \emph{Translation covariance.}
\item \emph{Locality for nested domains}: if $D'\subset D$ agrees with $D$
near regular marks $p,q$, the two initial laws agree until they encounter the
removed set.
\item \emph{Local target consistency}: for fixed $D,p$ and any two regular
targets $q_1,q_2$, the two initial laws agree when stopped on every
sufficiently small circle centered at $p$.
\item \emph{Geometric reversibility}: reversing an ordered path with law
$P_{D,p,q}$ gives law $P_{D,q,p}$.
\end{enumerate}
Then there is one fixed positive-definite tensor $Q_0$, up to scalar, with
$[Q_{D,p,q}]=[Q_0]$ for every admissible triple. Equivalently, one fixed
linear map transforms the entire family into ordinary chordal SLE$_6$.
No domain Markov hypothesis is needed for this conditional conclusion.
\end{theorem}

\begin{proof}
\emph{First, remove dependence on marks inside one domain.}
Fix $D$ and a regular start $p$. Local target consistency and local
identifiability make $[Q_{D,p,q}]$ independent of $q$.
Ordinary SLE$_6$ is geometrically reversible by \cite{ms}. Therefore reversal of (AA1)
retains its same tensor, and the assumed reversal of $P$ implies
\[
 [Q_{D,p,q}]=[Q_{D,q,p}].
\]
Given regular $p,p'$, choose a third regular point $q$. Changing the target
with fixed start $q$ and reversing at the two ends shows
\[
 [Q_{D,p,q}]=[Q_{D,q,p}]=[Q_{D,q,p'}]=[Q_{D,p',q}].
\]
Thus there is a tensor $[Q(D)]$ independent of both marks.

\emph{Second, compare nested domains with retained regular boundary arcs.}
If $D'\subset D$ agrees with $D$ near two distinct regular boundary points,
choose them as $p,q$. Locality makes the initial laws agree in some ball
about $p$, and local identifiability gives $[Q(D')]=[Q(D)]$.
The removed set may meet other parts of the boundary; it stays a positive
distance from the selected marks.

\emph{Third, create a horizontal boundary segment without changing the metric.}
For any $D\in\mathcal D$, choose a horizontal line $\{y=c\}$ with $\min_{\partial D}\operatorname{Im}z<c<\max_{\partial D}\operatorname{Im}z$, at a regular value of the boundary height and not through a corner, and intersect $D$ with
$\{y>c\}$. Such a regular value with nonempty intersection exists: on each nonhorizontal
analytic boundary arc the critical points are isolated, and horizontal arcs
contribute only finitely many excluded height values. Choose the level also
away from all corner heights. The analytic arcs are understood to extend
regularly to neighborhoods of their closed parameter intervals, so there
are only finitely many critical points unless an arc is horizontal.
There are finitely many transverse crossings, and a component $U$ of the
intersection is again a piecewise real-analytic Jordan domain. It contains an
original analytic boundary arc above the cutting line and also an open
horizontal segment of the cutting line on its boundary.
Choose two regular marks on the retained original arc. The second step gives
$[Q(U)]=[Q(D)]$.

Choose a sufficiently short subsegment of the horizontal boundary segment.
An upper half-disk $V$ on that diameter lies in $U$, agrees with $U$ near
two distinct interior diameter points, and belongs to $\mathcal D$.
The second step again gives $[Q(V)]=[Q(U)]$.

\emph{Finally, compare all horizontal upper half-disks.}
Translation covariance identifies their tensors for equal radii and arbitrary centers. Two upper half-disks with the same center and radii
$r<R$ are nested and agree near any two distinct points on the smaller open
diameter. The second step makes their tensors equal. Thus all such half-disks
have one tensor $[Q_0]$, and the third step shows $[Q(D)]=[Q_0]$ for every $D$.

Choose a fixed $A_0$ with $A_0^{\mathsf T}A_0=Q_0$. Equality of projective
tensors gives $A_{D,p,q}=c_{D,p,q}R_{D,p,q}A_0$ with $c_{D,p,q}>0$ and
$R_{D,p,q}$ orthogonal. Similarity covariance of SLE$_6$ in (AA1) removes
those extra factors, leaving exactly
$(A_0^{-1})_*\operatorname{SLE}_6(A_0D,A_0p,A_0q)$.
\end{proof}

\paragraph{Boundary regularity in the cutting step.}
A transverse horizontal cut of a piecewise real-analytic Jordan domain has finitely
many intersections with its boundary. Each component is simply connected:
a Jordan loop in the intersection has its interior in both the original
Jordan domain and the upper half-plane. Its boundary consists of finitely
many original arcs and cutting-line segments meeting at transverse corners,
so the component is Jordan with piecewise real-analytic boundary. At least one
retained original arc and one cutting segment occur in every nontrivial
component chosen in the proof. This also makes all marked-point comparisons
legitimate within the stated class.

\paragraph{Use for the percolation program.}
Local target consistency can be transferred from common-color discrete
couplings to simultaneous marginal curve limits without random-domain kernel
convergence, as established separately. Reversibility is a further assumption.
For a general lattice model with critical density different from $1/2$,
reversing the color-oriented exploration is not automatically the same as
swapping its start and target in a fixed convention: it also changes which
boundary arc has each color. No exact discrete color symmetry, nor limiting
reversal theorem, is supplied by the present argument.
The main unsolved analytic step remains establishing (AA1), or an alternative
local structure theorem, from the original axioms plus appropriate coherence.

% Research fragment; intended for inclusion in progress.tex.
\section{Continuous affine labels: deterministic support and source independence}
\label{sec:support-source}

\paragraph{Status.}
The result here is conditional on the deterministic affine-SLE ansatz and an
explicit continuity condition on its metric labels.  It does not apply to a
general exploration ensemble.  It gives a further restriction inside that
ansatz: a continuous conserved label cannot encode information in the source
prime end.  Source independence alone gives no target or domain independence; the final corollary adds explicit target-consistency and locality hypotheses to obtain a global metric.

\subsection{The precise continuity premise}

Let $T(D,p,q)$ be the determinant-one conformal metric tensor of the affine
label, as in (RA6), on all bounded simply connected domains with two distinct
prime ends.  We use the following sufficient premise, stated in a form that
is stable under the fixed affine change of coordinates needed for each
starting law.

\paragraph{Affine-stable marked-chart continuity.}
For every fixed invertible real-linear map $L$, suppose
$f_n:\mathbb H\to E_n$ are conformal maps onto bounded simply connected
domains, with a nondegenerate conformal limit
$f_n\to f:\mathbb H\to E$ locally uniformly in $\mathbb H$.  The marked
prime ends are $f_n(0),f_n(\infty)$ and $f(0),f(\infty)$, respectively;
these are prime-end extensions, whether or not they have Euclidean point
limits.  Require
\[
 T(LE_n,Lf_n(0),Lf_n(\infty))
   \longrightarrow T(LE,Lf(0),Lf(\infty)).
 \tag{SS1}
\]
Thus the topology is locally uniform convergence of inverse conformal
charts with the two designated boundary preimages $0,\infty$, and the
premise also applies after each fixed affine map.  This is an explicit
additional hypothesis; it is not inferred merely from DMP, locality, or
physical translation invariance.  An equivalent affine-stable marked kernel
notion may be used instead, provided it supplies (SS1).


\paragraph{Why fixed affine maps preserve this marked topology (proved using
standard Beltrami stability).}
In fact the displayed premise for all $L$ follows from its case $L=I$.
Here is the needed topological fact.  Write an orientation-preserving
real-linear map as $L(z)=\alpha z+\beta\bar z$, with
$|\beta/\alpha|<1$.  The Beltrami coefficients of $L\circ f_n$ are
\[
 \mu_n(z)=\frac{\beta}{\alpha}
             \frac{\overline{f_n'(z)}}{f_n'(z)}.
\]
Since the nonvanishing derivatives $f_n'$ converge locally uniformly to
$f'$, these coefficients converge pointwise in $\mathbb H$ to the analogous
$\mu$, all with the same bound $k<1$.  Extend them to the lower half-plane
by reflection, $\mu_n(\bar z)=\overline{\mu_n(z)}$.  Let $h_n$ be the
normalized quasiconformal sphere solutions fixing $0,1,\infty$.
The standard stability theorem for the measurable Riemann mapping theorem
(uniform bound $k<1$ and almost-everywhere convergence of coefficients)
gives locally uniform convergence of $h_n$ and their inverses to $h$ and
$h^{-1}$.  Reflection symmetry and normalized uniqueness imply that every
$h_n$ maps $\mathbb H$ onto itself and fixes the two marked boundary points
$0,\infty$.

The maps
\[
 \chi_n=L\circ f_n\circ h_n^{-1}:\mathbb H\longrightarrow LE_n
\]
are conformal by Beltrami factorization.  They converge locally uniformly
to the nondegenerate conformal map
$\chi=L\circ f\circ h^{-1}$, and their marked endpoints are precisely
$Lf_n(0),Lf_n(\infty)$.  This proves continuity of a fixed orientation-
preserving affine change of the physical domains in the marked conformal-
chart topology.  A reflection handles the other orientation.
Consequently ordinary marked-chart continuity of $T$ implies (SS1).
Keeping (SS1) explicit in the proof below makes clear that its charts are
holomorphic in the affine SLE coordinates, even when their physical pullbacks
are only quasiconformal.  The additional analytical input in this paragraph
is the standard normalized Beltrami stability theorem; no claim about
Euclidean uniform boundary convergence is needed.

\subsection{Full support transfers conservation to deterministic Loewner controls}

Fix a starting label $A=A(D,p,q)$ and a conformal map
$f:\mathbb H\to AD$ with $f(0)=Ap$, $f(\infty)=Aq$.  For any continuous
real function $w$ on $[0,t]$ with $w(0)=0$, let $g_t^w$ denote the associated
hydrodynamic Loewner map on its surviving half-plane domain.  Define the
marked target state by the inverse conformal chart
\[
 f_t^w(z)=f\bigl((g_t^w)^{-1}(z+w(t))\bigr),\qquad z\in\mathbb H.
 \tag{SS2}
\]
This state is a simply connected domain with two designated prime ends for
\emph{every} continuous $w$.  The assertion does not say that every such
$w$ generates a trace.  No trace is needed to define the domain state in
(SS2).  For the Brownian driving paths of SLE the state is its actual
remaining target domain.

The map $w\mapsto f_t^w$ is continuous from the uniform norm on $C[0,t]$
to locally uniform convergence of conformal maps on $\mathbb H$.
For example, the inverse map can be represented by the reverse ODE
\[
 \dot v_s=-\frac{2}{v_s-w(t-s)},\quad v_0=z,\quad
 (g_t^w)^{-1}(z)=v_t.
 \tag{SS3}
\]
Its imaginary part is nondecreasing.  On a compact set with
$\operatorname{Im}z\geq\delta>0$, denominators stay away from zero and the
usual ODE stability estimate gives locally uniform continuity in $w$.
Translations by $w(t)$ and composition with $f$ preserve this continuity.
The corresponding metric functional
\[
 \mathcal T_t(w)=
 T\bigl(A^{-1}f_t^w(\mathbb H),
         A^{-1}f_t^w(0),A^{-1}f_t^w(\infty)\bigr)
 \tag{SS4}
\]
is therefore continuous by (SS1), with $L=A^{-1}$.

\paragraph{Deterministic-support lemma (proved).}
If the affine-label ensemble has strong DMP and satisfies (SS1), then for
every $t>0$ and every continuous control $w\in C[0,t]$ with $w(0)=0$,
\[
 \mathcal T_t(w)=T(D,p,q).
 \tag{SS5}
\]
For the source-independence conclusion below it suffices to apply this to
smooth linear ramps.

\begin{proof}
Under the original affine SLE law the driving process in $AD$ is
$\sqrt6 B$.  The DMP conservation criterion (RA7) gives
$\mathcal T_t(\sqrt6 B)=T(D,p,q)$ almost surely.  Brownian measure has full
support on the space of continuous functions starting at zero in the
uniform topology.  If (SS5) failed at one control $w$, continuity of
$\mathcal T_t$ would give an open uniform neighborhood on which it differs
from the fixed tensor.  Brownian full support gives that neighborhood
positive probability, contradicting almost-sure conservation.
\end{proof}

One could avoid defining the functional at non-trace-generating controls
by restricting the statement to controls at which the state functional has
a continuous extension, and applying full support there.  In particular
smooth ramps suffice.  The all-prime-end ansatz already makes (SS4) defined
on every state, so the formulation above is simpler.

\subsection{Rapid driving ramps move the source at vanishing capacity}

\paragraph{Theorem (proved, within the stated ansatz).}
Under strong DMP and (SS1), the tensor $T(D,p,q)$ is independent of the
source prime end: for every bounded simply connected $D$ and every three
prime ends with $p,p'\ne q$,
\[
 T(D,p,q)=T(D,p',q).
 \tag{SS6}
\]
It may still depend on the target or on the domain; these are unresolved
by this argument.

\begin{proof}
Use the notation above and write $Ap'=f(a)$ for the unique finite
$a\in\mathbb R$ corresponding to the new source.  Choose $t_n\downarrow0$
and use the smooth control
\[
 w_n(s)=\frac{as}{t_n},\qquad 0\leq s\leq t_n.
\]
The deterministic-support lemma says that every resulting metric equals
$T(D,p,q)$.  From the reverse equation (SS3),
\[
 |(g_{t_n}^{w_n})^{-1}(z)-z|
 \leq\frac{2t_n}{\operatorname{Im}z}.
 \tag{SS7}
\]
Consequently, locally uniformly in $\mathbb H$,
\[
 f_{t_n}^{w_n}(z)
   = f\bigl((g_{t_n}^{w_n})^{-1}(z+a)\bigr)
   \longrightarrow f(z+a).
\]
The limiting chart has image $AD$, source $f(a)=Ap'$, and the same target
$f(\infty)=Aq$.  Applying (SS1) with $L=A^{-1}$ gives (SS6).
\end{proof}

The controls may have very large slopes as $t_n\to0$; this is harmless.
Each fixed ramp lies in the support of Brownian driving on its own finite
interval.  The argument uses the possibility of rare events, not their
uniform probability as $n\to\infty$.  The limiting source change is a
statement about marked prime-end charts, not Euclidean convergence of a
curve of vanishing diameter: the capacity vanishes while the tip may
travel far along a thin boundary region.

\subsection{A fixed metric in the affine ansatz without a reversal hypothesis}

\paragraph{Corollary (proved, conditional on the global affine ansatz).}
Suppose the ensemble has a deterministic affine-SLE label on \emph{all}
bounded simply connected domains and distinct prime ends, has strong DMP,
and its normalized tensor $T$ is continuous for marked inverse-conformal-
chart convergence.  Suppose in addition that its piecewise real-analytic
Jordan subfamily has translation covariance, locality for nested domains
retaining neighborhoods of regular marks, and local target consistency in
the sense used by the conditional global metric-rigidity theorem.  Then
there is one fixed determinant-one tensor $Q_0$ such that
\[
 T(D,p,q)=Q_0
\]
for every bounded simply connected marked domain.  No geometric reversal
hypothesis is needed for this conditional conclusion.

\begin{proof}
The source-independence theorem removes dependence on $p$.  On a regular
piecewise real-analytic Jordan domain, local target consistency together
with the local half-plane identifiability argument removes dependence on
$q$ as well.  Thus these domains have a tensor $Q(D)$ independent of both
marks.  The nested-cut argument in the preceding conditional global
metric-rigidity theorem now applies verbatim, without its reversal step:
nested locality identifies tensors of domains agreeing near two retained
regular marks; a transverse horizontal cut creates a horizontal segment;
a small upper half-disk on that segment retains the metric; translation
and nesting identify all such half-disks.  Hence all piecewise real-analytic
Jordan domains have the same tensor $Q_0$.

To extend to an arbitrary bounded simply connected $D$, choose a conformal
$f:\mathbb D\to D$ and distinct circle points $\zeta_p,\zeta_q$ representing
$p,q$.  For $r_n\uparrow1$, set
\[
 D_n=f(r_n\mathbb D),\qquad
 p_n=f(r_n\zeta_p),\quad q_n=f(r_n\zeta_q).
\]
Each $D_n$ is an analytic Jordan domain: $f$ extends holomorphically and
univalently across the compact circle $|z|=r_n$, with nonzero derivative.
Choose a conformal $m:\mathbb H\to\mathbb D$ with
$m(0)=\zeta_p$, $m(\infty)=\zeta_q$.  The marked inverse charts
\[
 f_n(z)=f(r_n m(z))
\]
converge locally uniformly to $f(m(z))$, giving exactly the desired marked
prime-end limit, even if $f$ has no Euclidean limit at one of the selected
circle points.  Tensor continuity and $T(D_n,p_n,q_n)=Q_0$ give
$T(D,p,q)=Q_0$.  The fixed-affine normalizer follows from SLE similarity
covariance as in the preceding theorem.
\end{proof}

This corollary removes a reversal assumption \emph{inside the global affine
ansatz}, at the price of strong DMP and a precise continuity premise.  It
still supplies no argument that the original three axioms, RSW, or local
target consistency force a marginal exploration law to belong to that ansatz.

\section{Unsuccessful routes, revised objectives, and next experiments}
\label{sec:next}

\paragraph{Why several tempting arguments fail.}
Independent site colors give exact discrete domain Markov and locality at
every occupation probability; these axioms do not by themselves encode
criticality. Near-critical scaling limits are therefore natural candidates,
but configuration convergence and a Markov theorem in the near-critical
parameter do not establish the required ordered exploration kernels at
random slit domains. We do not use this incomplete transfer as the main
counterexample. Nor does singularity relative to isotropic SLE alone exclude
an anisotropic affine image. A genuine dilation defect would do so.

A random mixture over fixed affine metrics is also not automatically domain
Markov: observing the initial curve can update the metric posterior. The
exact example avoids that failure through deterministic terminal information
that remains visible in every current domain. This should be distinguished
from unsupported claims that every stationary random environment yields a
domain-Markov annealed interface.

Applying Schramm's usual Brownian-driving argument directly would silently
assume that a conformally mapped slit domain has the same law as the original
domain. Physical translations do not imply that premise. The deletion-chart
identity above identifies the precise missing drift control. We have not
proved that the original axioms imply semimartingale driving, driving
martingales, or even a fixed underlying conformal structure.

\paragraph{A revised research question, not a proved classification.}
Restrict to nondegenerate continuous Loewner exploration laws with strong
domain Markov, fixed-target locality, physical translation covariance,
consistency under changing distant boundary conditions and target, and
continuity under controlled Jordan-domain approximation, including changes
near the terminal point. Require a genuine criticality condition, such as
nondegenerate crossing estimates uniform over all scales, formulated for a
common continuum coloring or a compatible exploration system. Does this
force one deterministic conformal metric, and then affine SLE$_6$?
Even the sufficiency of this strengthened collection is open here. One must
not promote its plausibility to a theorem. Full similarity covariance alone
is already defeated by the exact construction above.

\paragraph{Concrete priorities for the next session.}
The local common-coloring transfer is now proved in
Section~\ref{sec:local-target-transfer}. The next limit questions are
controlled marked-domain stability and conditional-kernel convergence in
random slit domains, ideally from uniform crossing estimates. The main
structural question is to derive affine-SLE marginals, or a different useful
local structure theorem, for arbitrary exploration laws. Once that ansatz
is available, the conditional common-metric theorem above removes the
domain dependence under explicit coherence and continuity hypotheses.
A parallel route is to seek an observable implying the driving-local-
martingale premise in Section~\ref{sec:loewner-locality}, possibly after
one affine normalization. That premise would finish identification by the
proved conditional theorem, but deriving it from coarse percolation
properties remains open. Finally examine whether later accessible frontier
germs can encode geometry despite target consistency: initial ordered-germ
decoding is now established, but the later unlabelled-boundary construction
is not. No full additional counterexample is claimed.

\paragraph{Reproducibility and preservation.}
The concise project handoff is \texttt{STATUS.md}. The polished counterexample
and global-deformation results are presented separately in \texttt{final.tex}.
Independent scratch investigations and compilation products are under
\texttt{work/}; this record contains the actual mathematical arguments and
does not require the conversation or those scratch files to be understood.
Successful LaTeX compilation checks presentation only. It is not evidence
that a mathematical proof is correct.


\begin{thebibliography}{99}
\bibitem{lawler} G. F. Lawler, \emph{Schramm--Loewner Evolution},
Park City lecture notes (2007), especially Lectures 3--4.
\url{https://math.uchicago.edu/~lawler/utah.pdf}.
\bibitem{yuan} Y. Yuan, \emph{Topological characterisations of Loewner traces},
Indiana Univ. Math. J. \textbf{71} (2022), 1027--1046,
Theorem 1.1 and Remark 1.4. \url{https://arxiv.org/abs/2003.05535}.
\bibitem{ms} J. Miller and S. Sheffield, \emph{Imaginary geometry III:
reversibility of SLE$_\kappa$ for $\kappa\in(4,8)$},
Ann. of Math. \textbf{184} (2016), 455--486.
\url{https://annals.math.princeton.edu/2016/184-2/p03}.
\bibitem{nw} P. Nolin and W. Werner,
\emph{Asymmetry of near-critical percolation interfaces},
J. Amer. Math. Soc. \textbf{22} (2009), 797--819.
\url{https://arxiv.org/abs/0710.1470}.
\bibitem{gps} C. Garban, G. Pete, and O. Schramm,
\emph{The scaling limits of near-critical and dynamical percolation},
J. Eur. Math. Soc. \textbf{20} (2018), 1195--1268,
especially Section 13.2. \url{https://arxiv.org/abs/1305.5526}.
\end{thebibliography}
\end{document}
